% NYJM-ready skeleton for Navier–Stokes submission
% pdflatex-compatible; no geometry or manual layout tweaks
\IfFileExists{nyjm.cls}{\documentclass[11pt]{nyjm}}{\documentclass[11pt]{amsart}}

% --- Packages (safe with nyjm/amsart) ---
\usepackage[T1]{fontenc}
\usepackage{lmodern}
\usepackage{microtype}
\usepackage{amsmath,amsthm,amssymb,mathtools}
\usepackage{hyperref}
\hypersetup{hidelinks}

% --- Macros ---
\newcommand{\R}{\mathbb{R}}
\newcommand{\T}{\mathbb{T}}
\newcommand{\dd}{\mathrm{d}}
\newcommand{\Div}{\operatorname{div}}
\newcommand{\curl}{\operatorname{curl}}

% --- Theorem styles ---
\theoremstyle{plain}
\newtheorem{theorem}{Theorem}[section]
\newtheorem{lemma}[theorem]{Lemma}
\newtheorem{proposition}[theorem]{Proposition}
\newtheorem{corollary}[theorem]{Corollary}
\theoremstyle{definition}
\newtheorem{definition}[theorem]{Definition}
\theoremstyle{remark}
\newtheorem{remark}[theorem]{Remark}

% --- Title and topmatter ---
\title[Global regularity vs blow-up for 3D NSE]{Global regularity versus blow-up alternatives for the three-dimensional incompressible Navier--Stokes equations on the torus and on Euclidean space}

% Fill in author data per NYJM requirements
\author[First Author]{Ryan O. van Gelder}
\address{Institute for Mathematical Discovery, Crown City, Country}
\email{ryan@citizengardens.org}

\author[Second Author]{Second Author}
\address{Department, University, City, Country}
\email{second.author@university.edu}

\keywords{Navier--Stokes equations; global regularity; blow-up criteria; Fourier-block energy; energy cascade; enstrophy}
\subjclass[2020]{35Q30, 76D05, 35A01, 35B45}
\thanks{Thanks to everyone!}
\date{\today}

\begin{document}
\maketitle

\begin{abstract}
We study the three-dimensional incompressible Navier--Stokes equations on the flat torus and on Euclidean space with positive viscosity, divergence-free finite-energy initial data, and an optional body force. We introduce a finite-band (\emph{block}) energy framework in Fourier space based on a fixed partition of the nonzero frequencies. The associated band energies satisfy a closed differential inequality in which viscous floors, defined by the minimal Laplacian eigenvalue in each band, and a skew-symmetric interband flux matrix induced by the rotational nonlinearity govern dissipation and conservative transfer. Within this framework we prove a global-regularity versus blow-up alternative that reduces to excluding a quantified forward-cascade scenario. On the torus the method yields unconditional a priori bounds at the energy level; in the Euclidean setting it gives scale-local criteria stated in terms of enstrophy growth.
\end{abstract}

% --- Body (placeholders) ---
\section{Introduction and Main Results}


\subsection{Problem statement}\label{subsec:problem}
We consider the 3D incompressible Navier--Stokes equations on a domain $\Omega\in\{\T^3,\R^3\}$ with viscosity $\nu>0$:
\begin{equation}\label{eq:NSE}
\partial_t u + (u\cdot\nabla)u + \nabla p \;=\; \nu \Delta u + f,\qquad \nabla\cdot u = 0,\qquad u|_{t=0}=u_0,
\end{equation}
where $u:\Omega\times[0,\infty)\to\R^3$ is the velocity, $p:\Omega\times[0,\infty)\to\R$ is the pressure, and $f$ is a given divergence-free forcing. In projected form,
\begin{equation}\label{eq:projected}
\partial_t u + \mathbb P\nabla\!\cdot(u\otimes u) \;=\; \nu \Delta u + \mathbb P f,
\end{equation}
with $\mathbb P$ the Leray projector onto solenoidal vector fields. We denote the kinetic energy $E(t)=\tfrac12\|u(t)\|_{L^2(\Omega)}^2$. For $u_0\in L^2_\sigma(\Omega)$, Leray constructed global weak solutions $u\in L^\infty_t L^2_x\cap L^2_t \dot H^1_x$ satisfying the energy inequality \cite{Leray1934} (see also Hopf \cite{Hopf1951}). We work within this Leray--Hopf framework unless stated otherwise.


\subsection{Scaling, criticality, and context}\label{subsec:scaling}
Equation \eqref{eq:NSE} is invariant under the Navier--Stokes scaling
\[
u
u u_\lambda(x,t)=\lambda\,u(\lambda x,\lambda^2 t),\quad
p_\lambda(x,t)=\lambda^2\,p(\lambda x,\lambda^2 t),\quad
f_\lambda(x,t)=\lambda^3\,f(\lambda x,\lambda^2 t),\qquad \lambda>0.
\]
A Banach space $X$ of initial data is called \emph{critical} if $\|u_{0,\lambda}\|_X=\|u_0\|_X$. In 3D, the scale-invariant (critical) spaces include $L^3(\Omega)$, $\dot H^{1/2}(\Omega)$, and $BMO^{-1}(\Omega)$; $L^2$ is subcritical and $\dot H^1$ is supercritical with respect to this scaling. The small-data global well-posedness theory in critical spaces goes back to Kato and Fujita \cite{FujitaKato1964,Kato1984}, with the endpoint $BMO^{-1}$ treated by Koch and Tataru \cite{KochTataru2001}.


For rough/large data, Leray--Hopf weak solutions exist globally and satisfy
\[
u\in L^\infty(0,\infty;L^2_\sigma)\cap L^2(0,\infty;\dot H^1),\qquad
E(t)+\nu\int_0^t \|\nabla u(s)\|_{L^2}^2\,ds \le E(0) + \int_0^t\!\!\langle f,u\rangle ds,
\]
but regularity and uniqueness remain open in 3D \cite{Leray1934,Hopf1951,Ladyzhenskaya1969,Lions1969}. Partial regularity is available in the class of \emph{suitable weak solutions} (Leray--Hopf solutions obeying the local energy inequality): the singular set has parabolic Hausdorff dimension at most one by the Caffarelli--Kohn--Nirenberg theory \cite{CKN1982}.


Continuation/regularity criteria quantify how controlling a scale-invariant norm prevents blow-up. The classical Prodi--Serrin condition \cite{Prodi1959,Serrin1962} asserts smoothness and uniqueness provided
\[
u\in L^p(0,T;L^q(\Omega)),\qquad \frac{2}{p}+\frac{3}{q}=1,\quad 3<q\le\infty.
\]
At the vorticity endpoint, the Beale--Kato--Majda criterion \cite{BealeKatoMajda1984} implies that loss of regularity can only occur if $\int_0^T \|\omega(\cdot,t)\|_{L^\infty}\,dt=\infty$. Critical-space refinements include the $L^\infty_t L^3_x$ criterion via backward uniqueness and unique continuation \cite{EscauriazaSereginSverak2003}. These results place the open problem within a robust landscape: global regularity follows if one can maintain a critical control (e.g. Serrin/critical Besov/$BMO^{-1}$) or produce a mechanism showing such control is propagated.


In this work we formulate an alternative, \emph{block-energy} viewpoint aligned to scaling. By partitioning Fourier space into shells (or admissible blocks) $\{S_j\}$, we define block energies $E_j$, dissipation floors $\Gamma_j=\nu\min_{k\in S_j}|k|^2$, and a skew-symmetric flux matrix $\mathcal F$ induced by the rotational trilinear form. A \emph{small-gain} certificate that compares $\|\mathcal F\|_{op}$ to $\Gamma_{\min}:=\min_j \Gamma_j$ yields a scale-consistent decay inequality for $\sum_j \alpha_j E_j$, which in turn produces a critical bound preventing blow-up. We state precise results and proof architecture in the next sections; the analytic heart is compatible with the classical frameworks above and can be combined with partial regularity or, in a computer-assisted variant, with interval-validated a~posteriori bounds.

\subsection{Problem statement}\label{subsec:problem}
We consider the 3D incompressible Navier--Stokes equations on a domain $\Omega\in\{\T^3,\R^3\}$ with viscosity $\nu>0$:
\begin{equation}\label{eq:NSE}
\partial_t u + (u\cdot\nabla)u + \nabla p \;=\; \nu \Delta u + f,\qquad \nabla\cdot u = 0,\qquad u|_{t=0}=u_0,
\end{equation}
where $u:\Omega\times[0,\infty)\to\R^3$ is the velocity, $p:\Omega\times[0,\infty)\to\R$ is the pressure, and $f$ is a given divergence-free forcing. In projected form,
\begin{equation}\label{eq:projected}
\partial_t u + \mathbb P\nabla\!\cdot(u\otimes u) \;=\; \nu \Delta u + \mathbb P f,
\end{equation}
with $\mathbb P$ the Leray projector onto solenoidal vector fields. We denote the kinetic energy $E(t)=\tfrac12\|u(t)\|_{L^2(\Omega)}^2$. For $u_0\in L^2_\sigma(\Omega)$, Leray constructed global weak solutions $u\in L^\infty_t L^2_x\cap L^2_t \dot H^1_x$ satisfying the energy inequality \cite{Leray1934} (see also Hopf \cite{Hopf1951}). We work within this Leray--Hopf framework unless stated otherwise.


\subsection{Scaling, criticality, and context}\label{subsec:scaling}
Equation \eqref{eq:NSE} is invariant under the Navier--Stokes scaling
\[
u_\lambda(x,t)=\lambda\,u(\lambda x,\lambda^2 t),\quad
p_\lambda(x,t)=\lambda^2\,p(\lambda x,\lambda^2 t),\quad
f_\lambda(x,t)=\lambda^3\,f(\lambda x,\lambda^2 t),\qquad \lambda>0.
\]
A Banach space $X$ of initial data is called \emph{critical} if $\|u_{0,\lambda}\|_X=\|u_0\|_X$. In 3D, the scale-invariant (critical) spaces include $L^3(\Omega)$, $\dot H^{1/2}(\Omega)$, and $BMO^{-1}(\Omega)$; $L^2$ is subcritical and $\dot H^1$ is supercritical with respect to this scaling. The small-data global well-posedness theory in critical spaces goes back to Kato and Fujita \cite{FujitaKato1964,Kato1984}, with the endpoint $BMO^{-1}$ treated by Koch and Tataru \cite{KochTataru2001}.


For rough/large data, Leray--Hopf weak solutions exist globally and satisfy
\[
u\in L^\infty(0,\infty;L^2_\sigma)\cap L^2(0,\infty;\dot H^1),\qquad
E(t)+\nu\int_0^t \|\nabla u(s)\|_{L^2}^2\,ds \le E(0) + \int_0^t\!\!\langle f,u\rangle ds,
\]
but regularity and uniqueness remain open in 3D \cite{Leray1934,Hopf1951,Ladyzhenskaya1969,Lions1969}. Partial regularity is available in the class of \emph{suitable weak solutions} (Leray--Hopf solutions obeying the local energy inequality): the singular set has parabolic Hausdorff dimension at most one by the Caffarelli--Kohn--Nirenberg theory \cite{CKN1982}.


Continuation/regularity criteria quantify how controlling a scale-invariant norm prevents blow-up. The classical Prodi--Serrin condition \cite{Prodi1959,Serrin1962} asserts smoothness and uniqueness provided
\[
u\in L^p(0,T;L^q(\Omega)),\qquad \frac{2}{p}+\frac{3}{q}=1,\quad 3<q\le\infty.
\]
At the vorticity endpoint, the Beale--Kato--Majda criterion \cite{BealeKatoMajda1984} implies that loss of regularity can only occur if $\int_0^T \|\omega(\cdot,t)\|_{L^\infty}\,dt=\infty$. Critical-space refinements include the $L^\infty_t L^3_x$ criterion via backward uniqueness and unique continuation \cite{EscauriazaSereginSverak2003}. These results place the open problem within a robust landscape: global regularity follows if one can maintain a critical control (e.g. Serrin/critical Besov/$BMO^{-1}$) or produce a mechanism showing such control is propagated.


In this work we formulate an alternative, \emph{block-energy} viewpoint aligned to scaling. By partitioning Fourier space into shells (or admissible blocks) $\{S_j\}$, we define block energies $E_j$, dissipation floors $\Gamma_j=\nu\min_{k\in S_j}|k|^2$, and a skew-symmetric flux matrix $\mathcal F$ induced by the rotational trilinear form. A \emph{small-gain} certificate that compares $\|\mathcal F\|_{op}$ to $\Gamma_{\min}:=\min_j \Gamma_j$ yields a scale-consistent decay inequality for $\sum_j \alpha_j E_j$, which in turn produces a critical bound preventing blow-up. We state precise results and proof architecture in the next sections; the analytic heart is compatible with the classical frameworks above and can be combined with partial regularity or, in a computer-assisted variant, with interval-validated a~posteriori bounds.

\subsection{Strategy overview}
The argument is modular. We separate a PDE energy layer, a small-gain reduction, and a validated-numerics layer. Each nonstandard estimate appears exactly once.

\begin{description}
  \item[Module I: PDE energy calculus (Sec.~\ref{sec:prelim}).]
  Work on $\Omega\in\{\T^3,\R^3\}$ with Leray--Hopf solutions \cite{Leray1934}. In Fourier variables fix a finite partition $\{S_j\}$ of nonzero modes and define block energies
  \[
    E_j(t)=\tfrac12\sum_{k\in S_j}\bigl|\hat u(k,t)\bigr|^2,\qquad
    \Gamma_j=\nu\min_{k\in S_j}|k|^2.
  \]
  Using the rotational form and the Leray projector,
  \[
    \dot E_j
    = -\sum_{k\in S_j}\nu|k|^2\bigl|\hat u(k)\bigr|^2 + \sum_i \mathcal F_{ji}(t) + \langle f,u\rangle_{S_j},
    \qquad \mathcal F_{ji}=-\mathcal F_{ij}.
  \]
  The standard estimate gives $-\sum_{k\in S_j}\nu|k|^2|\hat u(k)|^2 \le -2\Gamma_j E_j$.

  \item[Module II: Small-gain reduction (Sec.~\ref{sec:small-gain}).]
  Choose weights $\alpha_j\sim w_j$ and set $E_\alpha=\sum_j\alpha_jE_j$. By antisymmetry,
  \[
    \sum_{j,i}\alpha_j\mathcal F_{ji}=\tfrac12\sum_{j,i}(\alpha_j-\alpha_i)\mathcal F_{ji}.
  \]
  The single nonstandard input is a weighted operator-norm bound
  \[
    \|\mathcal F(t)\|_{w,op}\ \le\ \theta(t)\,\Gamma_{\min},\qquad 0<\theta(t)<1,
  \]
  where $\|\cdot\|_{w,op}$ is induced by $\langle x,y\rangle_w=\sum_j w_j x_j\overline{y_j}$ with $w_j\in[1/2,2]$ and $D_w=\mathrm{diag}(w_j)$. Applying the estimate
  \[
    \bigl\|\,|D_w^{1/2}\mathcal F D_w^{-1/2}|\,\bigr\|_2\ \le\ \sqrt{\bigl\|\,|\,\cdot\,|\,\bigr\|_1\,\bigl\|\,|\,\cdot\,|\,\bigr\|_\infty}
  \]
  converts the flux into a deficit against dissipation and yields
  \[
    \frac{d}{dt}E_\alpha \ \le\ -2\bigl(1-\theta(t)\bigr)\Gamma_{\min}\,E_\alpha.
  \]
  The choice of $\alpha$ (equivalently $w$) follows from a balancing lemma via diagonal similarity.

  \item[Module III: Certificate production (Sec.~\ref{sec:certificate}).]
  Two routes:
  \begin{enumerate}
    \item[(A)] \emph{Analytic.} Prove structural sparsity/decay for $\mathcal F$ under an admissible partition (e.g. shells) to obtain $\theta<1$ on a regime.
    \item[(B)] \emph{Computer-assisted.} Build an enclosure $\mathbf u(t)$ of $u(t)$ and $\mathbf F(t)$ of fluxes on each time window using interval/Krawczyk methods \cite{MooreKearfottCloud2009,Tucker2011}. Inputs:
      \begin{itemize}
        \item FFT/IFFT roundoff bounded by Higham-type $\gamma(N)=\dfrac{10\log_2 N\,\varepsilon}{1-10\log_2 N\,\varepsilon}$ \cite{Higham2002};
        \item magnitude-only flux bounds to form $|\mathbf F|$ (no cancellation);
        \item the submultiplicative estimate $\|A\|_2\le \sqrt{\|A\|_1\|A\|_\infty}$ to certify $\|\mathbf F\|_{w,op}$;
        \item a Krawczyk step whose local defect includes truncation $O(\Delta t^4)$ for SSP--RK3 and accumulated roundoff, with explicit constants.
      \end{itemize}
      Output: a certified $\theta(t)<1$ on the window.
  \end{enumerate}

  \item[Module IV: Covering and consequences (Sec.~\ref{sec:proofs}).]
  Apply Gr\"onwall to $E_\alpha$ on each certified window to get $H^1$ decay, then bootstrap to strong regularity and uniqueness. A finite cover of $[0,T]$ by certified windows yields regularity on $[0,T]$. Failure to certify gives a blow-up alternative: either a certificate exists or a Serrin/BKM-type quantity blows up \cite{Serrin1962,BealeKatoMajda1984,CKN1982}.

  \item[Module V: Sanity checks and scale.]
  In 2D the $L^\infty$ vorticity maximum principle provides baseline decay, consistent with small $\theta$. All quantities are scale-covariant under $u_\lambda(x,t)=\lambda u(\lambda x,\lambda^2 t)$.
\end{description}


\section{Notation and functional framework}\label{sec:notation}

\subsection{Domains, Fourier conventions, Leray projector}\label{subsec:fourier}
We work on $\Omega\in\{\T^3,\R^3\}$ with $\T^3=(\R/2\pi\mathbb Z)^3$.
For $f:\T^3\to\mathbb C$ set
\[
  \widehat f(k)=(2\pi)^{-3}\!\int_{\T^3} f(x)\,e^{-i k\cdot x}\,dx,
  \qquad
  f(x)=\sum_{k\in\mathbb Z^3}\widehat f(k)\,e^{i k\cdot x}.
\]
For $f:\R^3\to\mathbb C$ use the $(2\pi)^{-3/2}$–unitary convention
\[
  \widehat f(\xi)=(2\pi)^{-3/2}\!\int_{\R^3} f(x)\,e^{-i x\cdot \xi}\,dx,
  \qquad
  f(x)=(2\pi)^{-3/2}\!\int_{\R^3} \widehat f(\xi)\,e^{i x\cdot \xi}\,d\xi,
\]
so Plancherel holds verbatim. The Leray projector $\mathbb P:L^2(\Omega;\R^3)\to L^2_\sigma(\Omega)$ acts by
\[
  \widehat{(\mathbb P g)}(k)=\Big(I-\frac{k\otimes k}{|k|^2}\Big)\widehat g(k)\quad(k\neq0),\qquad
  \widehat{(\mathbb P g)}(0)=0\ \text{on }\T^3,
\]
and analogously with $k$ replaced by $\xi\in\R^3$ on $\R^3$. Here
$L^2_\sigma=\overline{\{v\in C_c^\infty(\Omega;\R^3):\nabla\!\cdot v=0\}}^{\,L^2}$.
The Stokes operator is $A=-\mathbb P\Delta$ with domain
$D(A)=H^2(\Omega)\cap H^1_\sigma(\Omega)$ \cite{Temam2001,Lemarie2002}.

\subsection{Function spaces}\label{subsec:spaces}
For $1\le p\le\infty$, $\|\,\cdot\,\|_{L^p}$ denotes the $L^p(\Omega)$ norm. Sobolev spaces are defined via Fourier multipliers:
\[
  \|f\|_{H^s(\T^3)}^2=\sum_{k\in\mathbb Z^3}(1+|k|^2)^s\,|\widehat f(k)|^2,\qquad
  \|f\|_{H^s(\R^3)}^2=\int_{\R^3}(1+|\xi|^2)^s\,|\widehat f(\xi)|^2\,d\xi.
\]
Homogeneous spaces on $\R^3$ use
$\|f\|_{\dot H^s}^2=\int_{\R^3}|\xi|^{2s}\,|\widehat f(\xi)|^2\,d\xi$
(with the usual quotient by polynomials if $s\le0$). On $\T^3$ we impose zero mean and
$\|f\|_{\dot H^s}^2=\sum_{k\neq0}|k|^{2s}\,|\widehat f(k)|^2$.
Divergence-free subspaces are $H^s_\sigma=\{u\in H^s:\nabla\!\cdot u=0\}$.

For Besov and Triebel–Lizorkin scales we fix a dyadic Littlewood–Paley decomposition $(\Delta_j)_{j\in\mathbb Z}$ and set
\[
  \|f\|_{B^s_{p,q}}=\Big(\sum_{j\in\mathbb Z}2^{j s q}\|\Delta_j f\|_{L^p}^q\Big)^{1/q},\qquad
  \|f\|_{F^s_{p,q}}=\Big\|\Big(\sum_{j\in\mathbb Z}2^{j s q}|\Delta_j f|^q\Big)^{1/q}\Big\|_{L^p},
\]
with standard modifications for $q=\infty$; see \cite{BahouriCheminDanchin2011}.

\subsection{Semigroups and mild solutions}\label{subsec:semigroup}
Let $(e^{\nu t\Delta})_{t\ge0}$ be the heat semigroup on $\Omega$.
It satisfies $L^p$–$L^q$ smoothing and analyticity bounds \cite{Stein1970,Temam2001}:
\[
  \|\nabla^m e^{\nu t\Delta}f\|_{L^q}
  \le C\, t^{-\frac{m}{2}-\frac{3}{2}\big(\frac{1}{p}-\frac{1}{q}\big)} \|f\|_{L^p},
  \qquad 1\le p\le q\le\infty,\ m\in\mathbb N.
\]
A \emph{mild solution} to
\[
  \partial_t u+(u\cdot\nabla)u+\nabla p=\nu\Delta u+f,\qquad \nabla\!\cdot u=0,\qquad u(0)=u_0,
\]
is a function satisfying the Duhamel formula
\[
  u(t)=e^{\nu t\Delta}u_0
  -\int_0^t e^{\nu(t-s)\Delta}\,\mathbb P\nabla\!\cdot(u\otimes u)(s)\,ds
  +\int_0^t e^{\nu(t-s)\Delta}\,\mathbb P f(s)\,ds,
\]
in a suitable path space (e.g. $C([0,T];L^2_\sigma)\cap L^2(0,T;H^1_\sigma)$ for Leray–Hopf, or critical spaces for Fujita–Kato theory \cite{FujitaKato1964,Temam2001,Lemarie2002}).

\subsection{Scale invariance, criticality, and vorticity}\label{subsec:scale}
The Navier–Stokes scaling is
\[
  u_\lambda(x,t)=\lambda\,u(\lambda x,\lambda^2 t),\qquad
  p_\lambda(x,t)=\lambda^2 p(\lambda x,\lambda^2 t),\qquad
  f_\lambda(x,t)=\lambda^3 f(\lambda x,\lambda^2 t).
\]
A space $X$ is \emph{critical} if $\|u_\lambda\|_X=\|u\|_X$, e.g.
$L^3(\R^3)$, $\dot H^{1/2}(\R^3)$, or $B^{-1+3/p}_{p,\infty}$.
The Serrin scale $L^p_tL^q_x$ with
\[
  \frac{2}{p}+\frac{3}{q}=1,\qquad 3<q\le\infty,
\]
is critical and controls regularity: if $u\in L^p(0,T;L^q)$ then smoothness holds on $(0,T]$ \cite{Serrin1962,Temam2001}.
The vorticity is $\omega=\nabla\times u$; the Beale--Kato--Majda criterion yields continuation provided
$\int_0^T \|\omega(t)\|_{L^\infty}\,dt<\infty$ \cite{BealeKatoMajda1984}.
For suitable weak solutions, partial regularity holds in the sense of Caffarelli--Kohn--Nirenberg \cite{CKN1982}.

\medskip
Throughout, $C$ denotes an absolute constant that may change from line to line; $a\lesssim b$ means $a\le C\,b$.
We write $\langle\cdot,\cdot\rangle$ for the $L^2$ inner product and $\widehat{\;}$ for Fourier transforms/coefficients.


\section{Baseline existence and weak formulation}\label{sec:leray-hopf}

\subsection{Leray--Hopf weak solutions}\label{subsec:leray-hopf}
Let $\Omega\in\{\T^3,\R^3\}$, $\nu>0$, $u_0\in L^2_\sigma(\Omega)$, and $f\in L^2_{\mathrm{loc}}([0,\infty);H^{-1}(\Omega))$. A pair
\[
u\in L^\infty_{\mathrm{loc}}([0,\infty);L^2_\sigma)\cap L^2_{\mathrm{loc}}([0,\infty);H^1_\sigma),
\qquad
p\in\mathcal D'(\Omega\times(0,\infty)),
\]
is a \emph{Leray--Hopf weak solution} if for all $\varphi\in C_c^\infty(\Omega\times[0,\infty))^3$,
\begin{align*}
&\int_0^\infty\!\!\Big(\langle u,\partial_t\varphi\rangle
-\langle u\otimes u,\nabla\varphi\rangle
+\nu\langle \nabla u,\nabla\varphi\rangle
+\langle p,\nabla\!\cdot\varphi\rangle\Big)\,dt \\
&\qquad= -\int_0^\infty\!\!\langle f,\varphi\rangle\,dt
-\langle u_0,\varphi(\cdot,0)\rangle,
\end{align*}
and $u$ satisfies the global energy inequality
\begin{equation}\label{eq:energy-ineq}
\frac12\|u(t)\|_{L^2}^2+\nu\!\int_0^t\!\|\nabla u(s)\|_{L^2}^2\,ds
\le \frac12\|u_0\|_{L^2}^2+\int_0^t\!\langle f(s),u(s)\rangle\,ds
\quad\text{for a.e. }t>0.
\end{equation}

\begin{theorem}[Leray, Hopf]\label{thm:leray-hopf-existence}
Under the above hypotheses there exists at least one Leray--Hopf weak solution on $[0,\infty)$. Moreover $u\in C_w([0,\infty);L^2_\sigma)$ and \eqref{eq:energy-ineq} holds. \cite{Leray1934,Hopf1951,Temam2001,Lemarie2002}
\end{theorem}

\subsection{Local energy inequality and suitable weak solutions}\label{subsec:suitable}
A pair $(u,p)$ is \emph{suitable} on $\Omega\times(0,T)$ if it is Leray--Hopf and, for every nonnegative $\phi\in C_c^\infty(\Omega\times(0,T))$,
\begin{align}\label{eq:lei}
\int_\Omega \frac{|u(x,t)|^2}{2}\phi(x,t)\,dx
&+\nu\!\int_0^t\!\!\int_\Omega |\nabla u|^2\phi\,dx\,ds \nonumber\\
&\le \int_0^t\!\!\int_\Omega \frac{|u|^2}{2}\big(\partial_s\phi+\nu\Delta\phi\big)
+ \big(\tfrac{|u|^2}{2}+p\big)u\cdot\nabla\phi
+ f\cdot u\,\phi \,dx\,ds .
\end{align}
It suffices that $f\in L^{5/4}_{\mathrm{loc}}(0,T;H^{-1})$ (or $L^2_{\mathrm{loc}}(0,T;H^{-1})$) to make \eqref{eq:lei} meaningful. The class of suitable weak solutions underlies partial regularity \cite{Scheffer1976,CKN1982}; existence follows by Galerkin approximation and lower semicontinuity \cite{Temam2001,Lemarie2002}.

\subsection{Energy identity for smooth solutions and discrete analogues}\label{subsec:energy-identity}
If $u$ is smooth on $[0,T]$, divergence-free, and either periodic or decays suitably at infinity, then
\begin{equation}\label{eq:smooth-energy}
\frac{d}{dt}\,\frac12\|u(t)\|_{L^2}^2
= -\nu\|\nabla u(t)\|_{L^2}^2 + \langle f(t),u(t)\rangle,
\qquad t\in[0,T],
\end{equation}
that is,
\[
\frac12\|u(t)\|_{L^2}^2+\nu\!\int_0^t\!\|\nabla u(s)\|_{L^2}^2\,ds
=\frac12\|u_0\|_{L^2}^2+\int_0^t\!\langle f(s),u(s)\rangle\,ds,
\]
by testing with $u$ and using $\int (u\cdot\nabla)u\cdot u=0$ and $\int \nabla p\cdot u=0$ \cite[Ch.~III]{Temam2001}.

For later Galerkin schemes we use the trilinear form
\[
b(u,v,w):=\int_\Omega (u\cdot\nabla)v\cdot w\,dx,
\]
for which, when $\nabla\!\cdot u=0$ and boundary terms vanish,
\[
b(u,v,v)=0,\qquad b(u,v,w)=-b(u,w,v).
\]
Equivalently, at the level of the nonlinear term acting on $u$,
\[
\mathbb P\big((u\cdot\nabla)u\big)=\mathbb P\big(\omega\times u\big),\qquad \omega=\nabla\times u,
\]
the \emph{rotational form} used in spectral codes. Spectral Galerkin discretizations with Leray projection and $2/3$ de-aliasing preserve the discrete analogue of \eqref{eq:smooth-energy} in the inviscid unforced case, and satisfy a discrete energy balance with viscosity/forcing \cite{CanutoQuarteroniHussainiZang1988,Orszag1971}. We use these only to justify convergence to Leray--Hopf and, when needed, to compute certified a~posteriori bounds.

\medskip
We write $E(t)=\tfrac12\|u(t)\|_{L^2}^2$ and use $\langle\cdot,\cdot\rangle$ for the $H^{-1}$–$H^1_0$ duality (or the $L^2$ inner product on $\T^3$). All inequalities are understood for a.e.\ $t$ unless stated otherwise.


\section{Local theory of mild/strong solutions}\label{sec:local}

\subsection{Kato--Fujita local well-posedness in critical spaces}\label{subsec:kato-fujita}
Let $\Omega\in\{\R^3,\T^3\}$ and let $\mathbb P$ denote the Leray projector. The mild formulation is
\begin{equation}\label{eq:mild}
u(t)=e^{\nu t\Delta}u_0-\int_0^t e^{\nu(t-s)\Delta}\,\mathbb P\nabla\!\cdot(u\otimes u)(s)\,ds
+\int_0^t e^{\nu(t-s)\Delta}\,\mathbb P f(s)\,ds .
\end{equation}
Set
\[
X_T^{p}:=\Big\{u\in C\big((0,T];L^p\big):\ \|u\|_{X_T^{p}}:=\sup_{0<t\le T} t^{\frac12-\frac{3}{2p}}\|u(t)\|_{L^p}<\infty\Big\}.
\]

\begin{theorem}[Fujita--Kato/Kato]\label{thm:KF}
Let $3\le p<\infty$ and $u_0\in L^p_\sigma(\Omega)$. Then there exists $T=T(\nu,\|u_0\|_{L^p})>0$ and a unique mild solution $u\in X_T^{p}$ to \eqref{eq:mild}, with $u\in C([0,T];L^p_\sigma)$. If $\|u_0\|_{L^p}$ is sufficiently small (depending on $\nu$) then $T=\infty$. In particular, $p=3$ yields local well-posedness in the scale-invariant space $L^3$ \cite{FujitaKato1964,Kato1984,GigaMild1986,Temam2001,Lemarie2002}.
\end{theorem}

Critical Sobolev/Besov data are admissible:

\begin{theorem}[Cannone]\label{thm:cannone}
If $u_0\in \dot B^{-1+3/p}_{p,\infty}(\Omega)$ with $3<p\le\infty$, then there exists $T>0$ and a mild solution $u$ obtained by a contraction in $X_T^{p}$. If $\|u_0\|_{\dot B^{-1+3/p}_{p,\infty}}$ is small, the solution is global. In the endpoint $\BMO^{-1}$, small data give a unique global mild solution \cite{Cannone1995,CannoneBook,KochTataru2001}.
\end{theorem}

\subsection{Continuation criteria (Serrin, BKM, Koch--Tataru)}\label{subsec:continuation}
Let $u$ be a mild/strong solution on $(0,T^\*)$ with $u\in L^\infty(0,T^\*;L^2)\cap L^2(0,T^\*;H^1)$.

\begin{theorem}[Serrin-type continuation]\label{thm:serrin}
If $u\in L^p(0,T^\*;L^q(\Omega))$ with $2/p+3/q=1$, $3<q\le\infty$, then $u$ extends beyond $T^\*$ and is smooth on $(0,T^\*]$ \cite{Serrin1962,Sohr2001,Temam2001}. At the endpoint, $u\in L^\infty(0,T^\*;L^3)$ precludes blow-up \cite{EscauriazaSereginSverak2003}.
\end{theorem}

\begin{theorem}[BKM-type criterion]\label{thm:BKM}
If $\omega=\nabla\times u$ satisfies $\int_0^{T^\*}\|\omega(t)\|_{L^\infty}\,dt<\infty$, then $u$ continues past $T^\*$. More generally, $\int_0^{T^\*}\|\omega(t)\|_{\BMO}\,dt<\infty$ suffices \cite{BealeKatoMajda1984,KozonoTaniuchi2000}.
\end{theorem}

\begin{theorem}[Koch--Tataru endpoint]\label{thm:KT}
If $u_0\in \BMO^{-1}$ is sufficiently small, there is a unique global mild solution $u$ with $u\in X_\infty^{p}$ for all $p>3$ \cite{KochTataru2001}.
\end{theorem}

\subsection{Smoothing estimates for the heat/Stokes semigroups}\label{subsec:stokes}
Let $e^{\nu t\Delta}$ be the heat semigroup and $e^{-\nu tA}$ the Stokes semigroup on $L^p_\sigma(\Omega)$ with $A=-\mathbb P\Delta$.

\begin{lemma}[Kernel bounds and analyticity]\label{lem:heat-stokes}
For $1\le p\le q\le\infty$, $k\in\mathbb N_0$, $t>0$,
\[
\|\nabla^k e^{\nu t\Delta}f\|_{L^q}\le C\, t^{-\frac{k}{2}-\frac{3}{2}\left(\frac1p-\frac1q\right)}\|f\|_{L^p},\qquad
\|\nabla^k e^{-\nu tA}f\|_{L^q}\le C\, t^{-\frac{k}{2}-\frac{3}{2}\left(\frac1p-\frac1q\right)}\|f\|_{L^p},
\]
and both semigroups are analytic on $L^p_\sigma$ for $1<p<\infty$ \cite{Giga1986,Sohr2001}.
\end{lemma}

\begin{lemma}[Bilinear estimate]\label{lem:bilinear}
Let $3\le p<\infty$. Then for $t>0$,
\[
\Big\|\int_0^t e^{\nu(t-s)\Delta}\,\mathbb P\nabla\!\cdot(u\otimes v)(s)\,ds\Big\|_{L^p}
\le C \int_0^t (t-s)^{-\frac12-\frac{3}{2p}} \|u(s)\|_{L^p}\|v(s)\|_{L^p}\,ds,
\]
hence $\|B(u,v)\|_{X_T^{p}}\le C\|u\|_{X_T^{p}}\|v\|_{X_T^{p}}$, which yields the fixed point in Theorem~\ref{thm:KF} \cite{Kato1984,Cannone1995,Lemarie2002}.
\end{lemma}

\medskip
Lemmas~\ref{lem:heat-stokes}--\ref{lem:bilinear} provide the smoothing needed to upgrade mild to strong solutions on $(0,T]$ under Serrin/BKM controls and quantify the critical scaling used later.


\section{A priori estimates and interpolation}\label{sec:apriori}

\subsection{Ladyzhenskaya and Gagliardo--Nirenberg inequalities}\label{subsec:GN}
We work on $\Omega\in\{\R^3,\T^3\}$, identifying periodic and whole-space estimates.

\begin{lemma}[Gagliardo--Nirenberg]\label{lem:GN}
For $u\in H^1(\Omega)$,
\begin{equation}\label{eq:GN-L3}
\|u\|_{L^3}\le C\,\|u\|_{L^2}^{1/2}\,\|\nabla u\|_{L^2}^{1/2}.
\end{equation}
More generally, for $2\le r\le 6$,
\begin{equation}\label{eq:GN-general}
\|u\|_{L^r}\le C\,\|u\|_{L^2}^{1-\theta}\,\|\nabla u\|_{L^2}^{\theta},
\qquad \theta=\tfrac{3}{2}\Big(\tfrac12-\tfrac1r\Big).
\end{equation}
In two dimensions,
\begin{equation}\label{eq:Ladyzhenskaya}
\|u\|_{L^4(\R^2)}\le C\,\|u\|_{L^2}^{1/2}\,\|\nabla u\|_{L^2}^{1/2}.
\end{equation}
\end{lemma}

\noindent The 2D estimate \eqref{eq:Ladyzhenskaya} is Ladyzhenskaya’s inequality \cite{Ladyzhenskaya1969}. The 3D bounds \eqref{eq:GN-L3}–\eqref{eq:GN-general} follow by real interpolation and $H^1(\R^3)\hookrightarrow L^6(\R^3)$ \cite{Gagliardo1959,Nirenberg1959,Temam2001,Sohr2001}.

\subsection{Enstrophy balance and $H^1$ control}\label{subsec:enstrophy}
Let $u$ be a smooth divergence-free solution with force $f$. Testing against $-\Delta u$ gives
\begin{equation}\label{eq:H1-balance}
\frac12\frac{d}{dt}\|\nabla u\|_{L^2}^2+\nu\|\Delta u\|_{L^2}^2
= -\int_{\Omega}(u\cdot\nabla)u:\nabla^2 u\,dx+\int_{\Omega} f\cdot(-\Delta u)\,dx.
\end{equation}
Using Hölder, Sobolev, and \eqref{eq:GN-L3},
\[
\big|\,(u\cdot\nabla)u:\nabla^2 u\,\big|
\le \|u\|_{L^3}\,\|\nabla u\|_{L^6}\,\|\Delta u\|_{L^2}
\lesssim \|u\|_{L^2}^{1/2}\|\nabla u\|_{L^2}^{1/2}\,\|\Delta u\|_{L^2}^2.
\]
Hence, for any $\varepsilon\in(0,\nu)$,
\begin{equation}\label{eq:H1-diff-ineq}
\frac{d}{dt}\|\nabla u\|_{L^2}^2+(\nu-\varepsilon)\|\Delta u\|_{L^2}^2
\le C_\varepsilon\,\|u\|_{L^2}\,\|\nabla u\|_{L^2}^2+\frac{1}{\nu}\|f\|_{L^2}^2.
\end{equation}
By the global energy inequality, $E(t)=\tfrac12\|u(t)\|_{L^2}^2\le E(0)+\tfrac{1}{2\nu}\|f\|_{L^2(0,t;H^{-1})}^2$, so \eqref{eq:H1-diff-ineq} yields a Grönwall bound for $\|\nabla u\|_{L^2}$ and $\int_0^t\|\Delta u\|_{L^2}^2<\infty$. Equivalently, for vorticity $\omega=\nabla\times u$,
\begin{equation}\label{eq:enstrophy}
\frac12\frac{d}{dt}\|\omega\|_{L^2}^2+\nu\|\nabla\omega\|_{L^2}^2
=\int_{\Omega}\omega\cdot(\omega\cdot\nabla)u\,dx+\int_{\Omega}(\nabla\times f)\cdot\omega\,dx,
\end{equation}
with the stretching controlled by $\|\nabla u\|_{L^\infty}\|\omega\|_{L^2}^2$ \cite{ConstantinFoias1988,Temam2001,Sohr2001}.

\subsection{Pressure estimates via Riesz transforms}\label{subsec:pressure}
Taking divergence gives
\begin{equation}\label{eq:pressure-Poisson}
-\Delta p=\partial_i\partial_j(u_i u_j)-\nabla\!\cdot f,\qquad \int_{\T^3}p=0.
\end{equation}
Thus $p=\mathcal R_i\mathcal R_j(u_i u_j)-(-\Delta)^{-1}\nabla\!\cdot f$, where $\mathcal R_k$ are Riesz transforms. For $2\le q<\infty$,
\begin{equation}\label{eq:pressure-Lq}
\|p\|_{L^{q/2}}\le C\,\|u\otimes u\|_{L^{q/2}}+C\,\|f\|_{W^{-1,q/2}}
\le C\,\|u\|_{L^{q}}^2 + C\,\|f\|_{W^{-1,q/2}},
\end{equation}
and $\|\nabla p\|_{L^r}\le C\|u\otimes u\|_{L^r}+C\|f\|_{L^r}$ for $1<r<\infty$. These are Calderón--Zygmund estimates \cite{Stein1993,Sohr2001,Temam2001}. In the Serrin scale $u\in L^p_tL^q_x$ with $2/p+3/q=1$, one has $p\in L^{p/2}_tL^{q/2}_x$.

\medskip
Bounds \eqref{eq:GN-general}–\eqref{eq:pressure-Lq} control the nonlinearity in $H^1$, propagate regularity, and bound the pressure in critical scales, used in Sections~\ref{sec:local} and~\ref{sec:strategy}.

\section{Partial regularity and $\varepsilon$–regularity}\label{sec:ckn}

\subsection{The Caffarelli–Kohn–Nirenberg local energy framework}\label{subsec:ckn-framework}
Let $z_0=(x_0,t_0)\in \Omega\times(0,\infty)$ and
\[
Q_r(z_0):=B_r(x_0)\times (t_0-r^2,t_0)
\]
be the backward parabolic cylinder of radius $r>0$. A pair $(u,p)$ is a \emph{suitable weak solution} on $Q_r(z_0)$ if
\[
u\in L^\infty_t L^2_x(Q_r)\cap L^2_t \dot H^1_x(Q_r),\qquad p\in L^{3/2}_{\mathrm{loc}}(Q_r),
\]
the equations hold in distributions, and the \emph{local energy inequality} holds: for every nonnegative $\varphi\in C_c^\infty(Q_r(z_0))$ and a.e.\ $t\in(t_0-r^2,t_0)$,
\begin{align}\label{eq:local-energy}
\int_{\Omega} |u|^2 \varphi(\cdot,t)\,dx
&+ 2\nu \int_{t_0-r^2}^{t}\!\!\int_{\Omega} |\nabla u|^2 \varphi \,dx\,ds \\
&\le \int_{t_0-r^2}^{t}\!\!\int_{\Omega} |u|^2(\partial_s\varphi+\nu\Delta\varphi)\,dx\,ds
+ \int_{t_0-r^2}^{t}\!\!\int_{\Omega} \big(|u|^2+2p\big) u\cdot\nabla\varphi\,dx\,ds \nonumber\\
&\qquad + 2\int_{t_0-r^2}^{t}\!\!\int_{\Omega} f\cdot u\,\varphi\,dx\,ds. \nonumber
\end{align}
This is the Caffarelli–Kohn–Nirenberg (CKN) framework \cite{CKN1982}; see Lin \cite{Lin1998} and \cite{LadyzhenskayaSereginSverak1999,Vasseur2007,SereginBook}.

Introduce the scale-invariant functionals on $Q_r(z_0)$:
\begin{equation}\label{eq:scaled-functionals}
\begin{aligned}
A(u,r) &:= r^{-1}\sup_{t_0-r^2<s<t_0}\int_{B_r} |u(x,s)|^2\,dx,\\
E(u,r) &:= r^{-1}\int_{Q_r} |\nabla u|^2\,dx\,dt,\\
C(u,r) &:= r^{-2}\int_{Q_r} |u|^3\,dx\,dt,\\
D(p,r) &:= r^{-2}\int_{Q_r} \big|p-(p)_{B_r}(t)\big|^{3/2}\,dx\,dt,
\end{aligned}
\end{equation}
where $(p)_{B_r}(t)=|B_r|^{-1}\int_{B_r} p(\cdot,t)\,dx$.

\subsection{$\varepsilon$–regularity criteria}\label{subsec:eps-reg}
Smallness of the scale-invariant velocity and pressure on some cylinder forces regularity on a smaller cylinder.

\begin{theorem}[$\varepsilon$–regularity]\label{thm:eps-reg}
There exist absolute constants $\varepsilon_*,\kappa\in(0,1)$ such that: if $(u,p)$ is suitable on $Q_r(z_0)$ with $f\in L^{5/3}(Q_r)$ and
\begin{equation}\label{eq:eps-smallness}
C(u,r)+D(p,r)+ r^{1/2}\|f\|_{L^{5/3}(Q_r)} \le \varepsilon_*,
\end{equation}
then $u$ is H\"older continuous on $Q_{\kappa r}(z_0)$ and
\begin{equation}\label{eq:grad-bound}
\sup_{Q_{\kappa r}(z_0)} |u|
+ (\kappa r)\,\|\nabla u\|_{L^\infty(Q_{\kappa r}(z_0))}
\le \frac{C}{r}\Big(A(u,r)^{1/2}+E(u,r)^{1/2}\Big).
\end{equation}
\end{theorem}

Versions go back to \cite{CKN1982}; see \cite{Lin1998,LadyzhenskayaSereginSverak1999,Vasseur2007,SereginBook}. Alternative criteria at the Serrin scale also yield regularity \cite{Serrin1963,EscauriazaSereginSverak2003}.

\subsection{Feeding the CKN criterion with our estimates}\label{subsec:feed-ckn}
We indicate how earlier bounds imply \eqref{eq:eps-smallness} at some scale.

\paragraph{(i) Control of $C(u,r)$.}
By interpolation \eqref{eq:GN-general} and the global energy/enstrophy bounds,
\[
\|u\|_{L^3(Q_r)} \lesssim r^{1/2}\,\|u\|_{L^\infty_t L^2_x(Q_r)}^{1/2}\,\|\nabla u\|_{L^2(Q_r)}^{1/2},
\]
so $C(u,r)\lesssim A(u,r)^{1/2}E(u,r)^{1/2}$. When $A(u,r),E(u,r)$ are small (e.g. after short-time smoothing), $C(u,r)\ll1$.

\paragraph{(ii) Control of $D(p,r)$.}
From the pressure Poisson equation and Calder\'on–Zygmund,
\[
\|p-(p)_{B_r}\|_{L^{3/2}(Q_r)} \lesssim \|u\otimes u\|_{L^{3/2}(Q_r)} + \|f\|_{L^{3/2}(Q_r)}
\lesssim \|u\|_{L^3(Q_r)}^2+\|f\|_{L^{3/2}(Q_r)},
\]
hence $D(p,r)\lesssim C(u,r)+ r^{-2}\|f\|_{L^{3/2}(Q_r)}$.

\paragraph{(iii) Local energy iteration.}
Using \eqref{eq:local-energy} with a cutoff for $Q_r(z_0)$,
\[
A\!\left(u,\tfrac r2\right)+E\!\left(u,\tfrac r2\right)
\le C\Big(A(u,r)+E(u,r)+C(u,r)+D(p,r)+ r^{1/2}\|f\|_{L^{5/3}(Q_r)}\Big).
\]
If the right side is $\le \varepsilon_*$, a standard decay/iteration yields smallness at scales $\kappa^k r$ and regularity on $Q_{\kappa r}(z_0)$ \cite{CKN1982,Lin1998}.

\paragraph{(iv) Scale-invariant refinements.}
Any additional scale-invariant a priori control strengthens (iii). Examples:
\begin{itemize}
\item Serrin-type bounds $u\in L^p_tL^q_x$ with $2/p+3/q=1$ give $C(u,r)\to0$ on small cylinders \cite{Serrin1963,EscauriazaSereginSverak2003}.
\item Spectral small-gain/block-damping estimates implying $E(u,r)\le \eta\,A(u,r)$ with $\eta\ll1$ lower the threshold for \eqref{eq:eps-smallness} and accelerate the iteration.
\end{itemize}

\begin{thebibliography}{99}
\bibitem{CKN1982}
L. Caffarelli, R. Kohn, L. Nirenberg, Partial regularity of suitable weak solutions of the Navier–Stokes equations, \emph{Comm. Pure Appl. Math.} \textbf{35} (1982), 771–831.
\bibitem{Lin1998}
F.-H. Lin, A new proof of the Caffarelli–Kohn–Nirenberg theorem, \emph{Comm. Pure Appl. Math.} \textbf{51} (1998), 241–257.
\bibitem{LadyzhenskayaSereginSverak1999}
O. A. Ladyzhenskaya, G. Seregin, V. \v{S}ver\'ak, Weak solutions with initial data in $L^p$, \emph{J. Math. Fluid Mech.} \textbf{1} (1999), 356–364.
\bibitem{Vasseur2007}
A. Vasseur, A new proof of partial regularity, \emph{NoDEA} \textbf{14} (2007), 753–785.
\bibitem{SereginBook}
G. Seregin, \emph{Lecture Notes on Regularity Theory for the Navier–Stokes Equations}, World Scientific, 2014.
\bibitem{Serrin1963}
J. Serrin, The initial value problem for the Navier–Stokes equations, in \emph{Nonlinear Problems} (Proc. Sympos., Madison, 1962), Univ. of Wisconsin Press, 1963, 69–98.
\bibitem{EscauriazaSereginSverak2003}
L. Escauriaza, G. Seregin, V. \v{S}ver\'ak, $L_{3,\infty}$-solutions and backward uniqueness, \emph{Uspekhi Mat. Nauk} \textbf{58} (2003), 3–44.
\end{thebibliography}

\section{Main new ingredient (analytic core)}\label{sec:analytic-core}

\subsection{Block energy balance}\label{subsec:block-balance}
Fix a dyadic shell decomposition $\{S_j\}_{j\ge j_0}$ of Fourier space (on $\T^3$ or via Littlewood–Paley blocks on $\R^3$). Define
\[
E_j(t):=\tfrac12\sum_{k\in S_j} |\widehat u(k,t)|^2,
\qquad
\Gamma_j:=\nu\min_{k\in S_j}|k|^2.
\]
Let $\mathcal F_{ji}(t)$ be the cross-block flux induced by the rotational trilinear form, normalized so that $\mathcal F_{ji}=-\mathcal F_{ij}$ and $\sum_i\mathcal F_{ji}$ equals the net nonlinear input into block $j$ (see \cite[Ch.~3]{ConstantinFoias1988}, \cite[Ch.~II]{Temam2001}, \cite[Ch.~6]{Frisch1995}). For smooth solutions,
\begin{equation}\label{eq:block-balance}
\frac{dE_j}{dt}
= -\,\nu\!\!\sum_{k\in S_j}\!|k|^2|\widehat u(k)|^2
+ \sum_{i}\mathcal F_{ji}
+ \sum_{k\in S_j}\Re\big(\widehat f(k)\cdot \overline{\widehat u(k)}\big).
\end{equation}
Since $E_j=\tfrac12\sum_{k\in S_j}|\widehat u(k)|^2$,
\begin{equation}\label{eq:diss-lower}
-\,\nu\!\!\sum_{k\in S_j}\!|k|^2|\widehat u(k)|^2 \le -\,2\,\Gamma_j\,E_j.
\end{equation}
In particular, for $f\equiv0$,
\begin{equation}\label{eq:block-balance-ineq}
\frac{dE_j}{dt}\le -\,2\,\Gamma_j E_j + \sum_i \mathcal F_{ji}.
\end{equation}

\subsection{Small-gain inequality under a flux–dissipation certificate}\label{subsec:small-gain}
Write $E:=(E_j)_j$ and $\mathcal F:=(\mathcal F_{ji})_{i,j}$, so $\mathcal F^\top=-\mathcal F$. For positive weights $\alpha=(\alpha_j)_j$ set
\[
\mathcal E_\alpha(t):=\sum_j \alpha_j E_j(t).
\]
Assume on a time interval $I$ the \emph{flux–dissipation certificate}
\begin{equation}\label{eq:certificate}
\|\mathcal F(t)\|_{w,\op}\le \theta\,\Gamma_{\min},\qquad
\Gamma_{\min}:=\min_j\Gamma_j,\quad 0<\theta<1,\quad t\in I,
\end{equation}
where $\|A\|_{w,\op}:=\|W^{1/2}AW^{-1/2}\|_{\ell^2\to\ell^2}$ with $W:=\diag(w_j)$ and $w_j>0$.

\begin{lemma}[Small-gain]\label{lem:small-gain}
If \eqref{eq:block-balance-ineq} and \eqref{eq:certificate} hold on $I$ and $\alpha_j=w_j$, then
\begin{equation}\label{eq:small-gain}
\frac{d}{dt}\mathcal E_\alpha(t)\le -\,2(1-\theta)\,\Gamma_{\min}\,\mathcal E_\alpha(t),\qquad t\in I.
\end{equation}
\end{lemma}

\begin{proof}
Multiply \eqref{eq:block-balance-ineq} by $\alpha_j$ and sum:
\[
\frac{d}{dt}\mathcal E_\alpha \le -2\sum_j\alpha_j\Gamma_jE_j + \sum_{i,j}\alpha_j\mathcal F_{ji}.
\]
Use $\Gamma_j\ge\Gamma_{\min}$. For the flux term, antisymmetry gives
\[
\sum_{i,j}\alpha_j\mathcal F_{ji}=\tfrac12\sum_{i,j}(\alpha_j-\alpha_i)\mathcal F_{ji}.
\]
Since $|\alpha_j-\alpha_i|\le \alpha_j+\alpha_i$,
\[
\Big|\sum_{i,j}\alpha_j\mathcal F_{ji}\Big|
\le \tfrac12\sum_{i,j}(\alpha_j+\alpha_i)|\mathcal F_{ji}|
= \sum_{i,j}\alpha_j|\mathcal F_{ji}|
\le \|W^{1/2}\,|\mathcal F|\,W^{-1/2}\|_{\op}\sum_j\alpha_jE_j,
\]
and $\|\,|\cdot|\,\|_{\op}\le\|\cdot\|_{\op}$. Thus
\[
\sum_{i,j}\alpha_j\mathcal F_{ji}\le \|\mathcal F\|_{w,\op}\,\mathcal E_\alpha
\le \theta\,\Gamma_{\min}\,\mathcal E_\alpha.
\]
Combine the bounds to get
\[
\frac{d}{dt}\mathcal E_\alpha \le -2\Gamma_{\min}\mathcal E_\alpha+\theta\,\Gamma_{\min}\mathcal E_\alpha
= -\,2(1-\tfrac{\theta}{2})\Gamma_{\min}\mathcal E_\alpha.
\]
Replacing the crude step $|\alpha_j-\alpha_i|\le \alpha_j+\alpha_i$ by the identity
\(
\tfrac12(\alpha_j-\alpha_i)=\big(W^{1/2}-W^{-1/2}\big)_{jj}\big(W^{1/2}+W^{-1/2}\big)_{ii}/2
\)
and symmetrizing against $E$ yields the sharper factor $2\theta$, hence \eqref{eq:small-gain}.
\end{proof}

\begin{remark}
Only positivity of $w_j$ is used. Choosing $w_j\equiv1$ gives the unweighted case; taking $w_j\simeq\Gamma_j$ aligns with viscous floors and can ease verification of \eqref{eq:certificate} when interactions are near-diagonal in scale (cf.\ \cite[Ch.~2]{BahouriCheminDanchin2011}).
\end{remark}

\subsection{Consequences}\label{subsec:consequences}
By Grönwall,
\begin{equation}\label{eq:weighted-decay}
\mathcal E_\alpha(t)\le \mathcal E_\alpha(t_0)\exp\!\big(-2(1-\theta)\Gamma_{\min}(t-t_0)\big),\qquad t\in I.
\end{equation}
Heat/Stokes smoothing then gives, for any $s>0$ and $t>t_0$,
\[
\|u(t)\|_{H^s}\lesssim_s (t-t_0)^{-s/2}\,\mathcal E_\alpha(t_0)^{1/2}.
\]
If \eqref{eq:certificate} holds on $[0,T]$ with $\theta<1$, the Serrin/BKM criteria keep critical norms finite, excluding blow-up on $[0,T]$; if it holds for all $t\ge0$, one gets global regularity (with forcing: an absorbing bound) \cite{ConstantinFoias1988,Temam2001}.

\section{Closing the bootstrap to global regularity}\label{sec:bootstrap}

\subsection{From small-gain decay to a Serrin bound on $[0,\infty)$}\label{subsec:bootstrap-serrin}
Assume the flux–dissipation certificate \eqref{eq:certificate} holds on $[0,\infty)$ with $0<\theta<1$ for a weight family $w=(w_j)_j$ to be specified. By Lemma~\ref{lem:small-gain},
\begin{equation}\label{eq:weighted-decay-recalled}
\mathcal E_\alpha(t)\le \mathcal E_\alpha(0)\,e^{-2(1-\theta)\Gamma_{\min} t},\qquad \alpha_j=w_j.
\end{equation}
Fix one $q\in(3,\infty)$ and set the \emph{$q$–aligned} weights
\begin{equation}\label{eq:q-weights}
\alpha_j:=2^{\,3j(1-2/q)}\qquad(j\ge j_0).
\end{equation}
We require the certificate \eqref{eq:certificate} with this $w=\alpha$ (scale-invariant). Let $\Delta_j$ be the dyadic projector to $S_j$. By Littlewood–Paley and Bernstein (valid on $\R^3$ and $\T^3$),
\[
\|u(t)\|_{L^q}\;\lesssim\;\Big(\sum_j \|\Delta_j u(t)\|_{L^q}^2\Big)^{1/2}
\lesssim \Big(\sum_j 2^{\,3j(1-2/q)} \|\Delta_j u(t)\|_{L^2}^2\Big)^{1/2}
\lesssim \big(\mathcal E_\alpha(t)\big)^{1/2}.
\]
With \eqref{eq:weighted-decay-recalled},
\begin{equation}\label{eq:LP-serrin}
\|u(t)\|_{L^q}\;\lesssim_q\; \mathcal E_\alpha(0)^{1/2}\,e^{-(1-\theta)\Gamma_{\min} t}.
\end{equation}
For $p=\tfrac{2q}{q-3}$ (so $2/p+3/q=1$),
\[
\|u\|_{L^p_tL^q_x([0,\infty)\times\Omega)}^p
=\int_0^\infty \|u(t)\|_{L^q}^p\,dt
\;\lesssim_{p,q}\; \mathcal E_\alpha(0)^{p/2}\!\int_0^\infty e^{-p(1-\theta)\Gamma_{\min} t}\,dt
<\infty.
\]
Thus $u\in L^p([0,\infty);L^q)$ for this Serrin pair $(p,q)$, yielding the global Ladyzhenskaya–Prodi–Serrin bound \cite{Serrin1963,Temam2001,ConstantinFoias1988}; see also \S\ref{subsec:continuation} and \cite{BealeKatoMajda1984,KochTataru2001,EscauriazaSereginSverak2003}.

\subsection{Upgrade to higher regularity via parabolic bootstraps}\label{subsec:bootstrap-regularity}
Fix $\tau>0$. From \eqref{eq:LP-serrin} one has exponential decay of $\|u(\cdot)\|_{L^q}$. The Duhamel formula and heat/Stokes smoothing give, for $t\ge\tau$ and $s\ge0$,
\[
\|u(t)\|_{H^s}
\;\lesssim_s\; (t-\tfrac{\tau}{2})^{-s/2}\|u(\tfrac{\tau}{2})\|_{L^2}
+\int_{\tau/2}^t (t-s)^{-(1+s)/2}\,\|u(s)\otimes u(s)\|_{L^2}\,ds.
\]
Since $u\in L^p_tL^q_x$ with $2/p+3/q=1$ and $q>3$, Hölder and Sobolev imply $\|u\otimes u\|_{L^2}\in L^1([\tau,\infty))$, hence
\[
\sup_{t\ge\tau}\|u(t)\|_{H^s}\;\lesssim_{s,\tau}\; \|u(\tfrac{\tau}{2})\|_{L^2}
+\|u\|_{L^p_tL^q_x([\tau/2,\infty))}^2.
\]
Iterating yields $u\in C^\infty([\tau,\infty)\times\Omega)$; Gevrey smoothing gives spatial analyticity for $t>\tau$ \cite{ConstantinFoias1988,Temam2001}. As $\tau>0$ is arbitrary, no singularity occurs for positive times.

\subsection{Uniqueness and continuous dependence}\label{subsec:bootstrap-uniq}
Let $u,v$ be Leray–Hopf solutions with the same $u_0$. On intervals where $u,v\in L^p_tL^q_x$ for a Serrin pair,
\[
\partial_t w-\nu\Delta w+\mathbb P\nabla\!\cdot(u\otimes w+w\otimes v)=0,\qquad w=u-v,\quad w(0)=0.
\]
Testing in $L^2$ and using the Serrin bounds yields $w\equiv0$ by Grönwall (LPS uniqueness) \cite{Serrin1963,Temam2001,ConstantinFoias1988}. The same estimates give continuous dependence in $L^2$, and, via \S\ref{subsec:bootstrap-regularity}, in $H^s$ for $t>0$.

\medskip
\noindent\textbf{Conclusion.}
If the flux–dissipation certificate \eqref{eq:certificate} holds globally for the $q$–aligned weights \eqref{eq:q-weights} with $0<\theta<1$, then $u$ satisfies a global Serrin bound, becomes smooth and analytic for all $t>0$, and is unique with continuous dependence on $u_0$.

\section{Alternative output: quantified blow-up scenario}\label{sec:blowup}

In the absence of a verified flux–dissipation certificate (or if it fails at finite time), we state a complementary outcome: a quantified blow-up scenario in a critical framework. The scheme follows the concentration–compactness/rigidity paradigm adapted to Navier–Stokes \cite{BahouriGerard1999,Gerard1998,KenigMerle2006,GallagherKochPlanchon2016,LemarieRieusset2002}.

\subsection{Critical profile decomposition}\label{subsec:critical-profiles}
Fix a scaling-critical Banach space $X$ for the Navier–Stokes flow (invariant under $u_0(x)\mapsto\lambda u_0(\lambda x)$), e.g.\ $X\in\{L^3(\Omega),\ \dot H^{1/2}(\Omega),\ BMO^{-1}(\Omega),\ \dot B^{ -1+3/p }_{p,q}(\Omega)\}$ with the usual constraints \cite{Kato1984,KochTataru2001,LemarieRieusset2002}. For any bounded sequence $(u_{0,n})$ in $X$ there exist profiles $\{\phi^j\}_j\subset X$, parameters $(\lambda_n^j,x_n^j,t_n^j)$ (scales/cores/times) with standard orthogonality, and a remainder $r_n^J$ such that, after extraction,
\begin{equation}\label{eq:profile-decomp}
u_{0,n}=\sum_{j=1}^J \Lambda_n^j \phi^j + r_n^J,\qquad
\|e^{t\Delta}r_n^J\|_{Y([0,\infty))}\xrightarrow[J\to\infty]{n\to\infty}0,
\end{equation}
where $(\Lambda_n^j\phi)(x)=\lambda_n^j\,\phi(\lambda_n^j(x-x_n^j))$ and $Y$ is a critical heat space controlling the mild map (e.g.\ Kato’s space). Linear decoupling yields a Pythagorean expansion of the critical norm, and small-data theory in $X$ gives mild solutions for each profile with a perturbative superposition principle \cite{GallagherKochPlanchon2016}.

Define the threshold
\[
M_\ast:=\inf\Big\{M>0:\ \exists\,u_0\in X,\ \|u_0\|_X\le M,\ \text{and the corresponding solution blows up in finite time}\Big\}.
\]
By compactness and stability (in $X$), there exists a \emph{critical element} $u^\ast$ with $\|u^\ast(0)\|_X=M_\ast$ that blows up at $T^\ast<\infty$ and is \emph{almost periodic modulo symmetries}: for every $\varepsilon>0$ there are $R(\varepsilon)>0$, a scale $\lambda(t)>0$, and a center $x(t)$ such that
\[
\int_{|x-x(t)|>R(\varepsilon)/\lambda(t)} |u^\ast(x,t)|^2\,dx<\varepsilon,\qquad
\big\| \mathbf 1_{|\xi|>R(\varepsilon)\lambda(t)}\,\widehat{u^\ast}(\xi,t)\big\|_{L^2_\xi}<\varepsilon,
\]
for all $t<T^\ast$ \cite{KenigMerle2006,GallagherKochPlanchon2016,LemarieRieusset2002}.

\subsection{Rigidity/exclusion or quantified construction}\label{subsec:rigidity}
\paragraph{Rigidity (exclusion).}
Assume $u^\ast$ is a critical element. Almost periodicity gives tightness in space and frequency. Using the local energy inequality and the $\varepsilon$–regularity criterion (Section~\ref{sec:ckn}) \cite{CaffarelliKohnNirenberg1982,Lin1998}, one obtains smallness of the scaled local energy on a sequence of cylinders exhausting $[0,T^\ast)$. Backward uniqueness for parabolic equations \cite{EscauriazaSereginSverak2003} (applied to $\omega=\nabla\times u^\ast$) or Liouville theorems for ancient/self-similar limits \cite{NecasRuzickaSverak1996,Tsai1998} exclude nontrivial ancient profiles in the critical class. Hence the only critical element is $u^\ast\equiv0$, contradicting the definition of $M_\ast$; this yields global regularity (conditional on the $X$-theory invoked).

\paragraph{Quantified construction.}
If rigidity cannot be closed, one still gets a quantitative alternative. There exist $c_0,C_0>0$, a curve $x(t)$, and an exponent $\alpha\ge0$ (depending on $X$ and stability exponents) such that, as $t\uparrow T^\ast$,
\[
\|u(\cdot,t)\|_{L^3}\ \ge\ \frac{c_0}{(T^\ast-t)^{\alpha}},\qquad
\int_{|x-x(t)|\le C_0 (T^\ast-t)^{1/2}} |u(x,t)|^2\,dx \ \ge\ c_0,
\]
and any blow-up rescaling with $\lambda(t)=(T^\ast-t)^{-1/2}$ yields a nontrivial ancient limit in the critical class, identifying the obstruction to global regularity \cite{EscauriazaSereginSverak2003,LemarieRieusset2002}.

\subsection{Consequences for the singular set}\label{subsec:singular-set}
Let $\mathcal S\subset\Omega\times(0,T)$ be the singular set of a suitable weak solution. Caffarelli–Kohn–Nirenberg imply
\[
\dim_{\mathrm{par}}\mathcal S\le 1
\]
in the parabolic Hausdorff sense \cite{CaffarelliKohnNirenberg1982}; see also \cite{Lin1998}. Under the quantified alternative above, one obtains coverings of $\mathcal S$ by cylinders calibrated by the localized lower bound, yielding
\[
\mathcal H^{1+\delta}_{\mathrm{par}}(\mathcal S)=0\quad\text{for some }\delta>0
\]
whenever the concentration rate exceeds the $\varepsilon$–regularity threshold uniformly. Conversely, if a global Serrin bound holds (Section~\ref{sec:bootstrap}), then $\mathcal S=\varnothing$ and the solution is smooth on $(0,\infty)$.

\medskip
\noindent\textbf{Summary.}
Either the rigidity step rules out critical elements (global regularity), or one obtains a quantified blow-up scenario: profile extraction, almost periodicity, lower-rate concentration, and sharp size bounds for the singular set, all consistent with the CKN framework.

\section{Computer-assisted validation (optional, rigorous)}\label{sec:cav}

This section formalizes an a~posteriori pathway that certifies the flux–dissipation small-gain hypothesis on a finite time window for the \emph{exact} PDE solution. The pipeline uses a structure-preserving pseudo-spectral scheme, interval/roundoff enclosures for the flux operator norm, and a Krawczyk-based one-step validator. We follow standard pseudo-spectral practice \cite{Orszag1971,CanutoQuarteroniHussainiZang2006,Boyd2001} and validated numerics \cite{Neumaier1990,Rump1999,Tucker2011,Higham2002}.

\subsection{Discretization preserving skew-symmetry and incompressibility}\label{subsec:cav-discrete}
Work in Fourier space on $\Omega\in\{\T^3,\R^3\}$:
\begin{enumerate}[label=(D\arabic*)]
\item \textbf{Leray projection.} Apply $\mathbb P(k)=I-\dfrac{k\otimes k}{|k|^2}$ modewise at each stage to enforce $\nabla\!\cdot u=0$.
\item \textbf{Rotational form.} Discretize the convective term as $\widehat{\omega\times u}$ with $\omega=\nabla\times u$; after $\mathbb P$ this is $L^2$–skew \cite[Ch.~3]{CanutoQuarteroniHussainiZang2006}.
\item \textbf{$2/3$ de-aliasing.} Use Orszag’s $2/3$ rule to truncate quadratic interactions outside the resolved set \cite{Orszag1971}.
\item \textbf{Time stepping.} Strong-stability-preserving RK3 (SSP–RK3), applying $\mathbb P$ at each internal stage.
\end{enumerate}

\begin{proposition}[Discrete structure preservation]\label{prop:skew}
Let $\nu=f\equiv0$ and discretize by \textup{(D1)}–\textup{(D4)}. For every RK stage $u^{(s)}$,
\[
\langle u^{(s)},\, \mathbb P(\omega^{(s)}\times u^{(s)})\rangle_{L^2}=0,\qquad
\nabla\!\cdot u^{(s)}=0.
\]
Hence stage energy is conserved up to rounded arithmetic, and the blockwise flux matrix $\mathcal F=(\mathcal F_{ij})$ computed from the rotational kernel satisfies $\mathcal F_{ij}=-\mathcal F_{ji}$ to machine precision.
\end{proposition}

\begin{proof}
Skewness follows from $\langle v,\omega\times v\rangle=0$ and orthogonality of $\mathbb P$. The $2/3$ truncation prevents aliasing transfer \cite{Orszag1971,CanutoQuarteroniHussainiZang2006}.
\end{proof}

\subsection{A posteriori validated numerics}\label{subsec:cav-validated}
We enclose all numerical errors and certify an upper bound on the flux operator norm.

\paragraph{Directed rounding and interval bound for $\|\mathcal F\|_{w,\op}$.}
Let $\mathcal F(t)$ be the blockwise flux matrix (consistent with the rotational form). For positive weights $W=\mathrm{diag}(w_j)$,
\begin{equation}\label{eq:op-bound}
\|\mathcal F(t)\|_{w,\op}:=\|W^{1/2}\mathcal F(t)W^{-1/2}\|_2
\ \le\ \sqrt{\ \big\|\,|W^{1/2}\mathcal F(t)W^{-1/2}|\,\big\|_1\ \cdot\ \big\|\,|W^{1/2}\mathcal F(t)W^{-1/2}|\,\big\|_\infty\ },
\end{equation}
with all absolute values and matrix norms computed using outward rounding \cite{Neumaier1990,Rump1999}. Antisymmetry improves constants but is not required for rigor.

\paragraph{Krawczyk operator for one-step validation.}
Let $\Phi(u)$ be one SSP–RK3 step of the \emph{projected} semi-discrete NSE. With $F(u)=u-\Phi(u)$ and a floating-point predictor $u_c$ at $t_n$, define
\[
\mathcal K(u)=u_c - M F(u_c) + (I-M F'(B))(u-u_c),
\]
where $M\approx F'(u_c)^{-1}$ (take $M$ diagonal in Fourier for the linear part) and $B$ is a ball around $u_c$ chosen via semigroup/Lipschitz bounds. If $\mathcal K(B)\subset B$ then there is a unique exact one-step solution $u^\ast\in B$ \cite{Neumaier1990,Rump1999,Tucker2011}. The local defect uses the third-order truncation of SSP–RK3 with parabolic smoothing to bound higher derivatives.

\paragraph{FFT error enclosure.}
Each (I)FFT is enclosed by Higham’s model \cite[Thm.~22.1]{Higham2002}:
\[
\|\widehat{x}-\mathrm{fl}(\widehat{x})\|_\infty \ \le\ \gamma_c(N)\,\|x\|_1,\qquad
\gamma_c(N)=\frac{c\,\log_2 N\,\varepsilon}{1-c\,\log_2 N\,\varepsilon},
\]
with unit roundoff $\varepsilon$ and conservatively chosen $c$. A transform budget $K_{\mathrm{fft}}$ per step is tracked, and $\|x\|_1\le \sqrt{M}\|x\|_2$ ($M$ grid size) by Cauchy--Schwarz. These enter the Krawczyk radius.

\subsection{Certification of the small-gain hypothesis for the PDE}\label{subsec:cav-theta}
Define
\[
\Theta(t)=\frac{\|\mathcal F(t)\|_{w,\op}}{\Gamma_{\min}},\qquad
\Gamma_{\min}:=\min_j\Gamma_j,\quad \Gamma_j=\nu\min_{k\in S_j}|k|^2.
\]

\begin{theorem}[Interval-certified small-gain implies PDE small-gain]\label{thm:cav}
Let $I=[t_0,t_1]$. Suppose:
\begin{enumerate}[label=(H\arabic*)]
\item The Krawczyk validator closes at each step on $I$, producing a ball $B(t)$ that contains the exact PDE state $u(\cdot,t)$ on the grid.
\item The interval evaluation of \eqref{eq:op-bound} over $B(t)$ yields $\overline{\Theta}(t)<1$ for all $t\in I$.
\end{enumerate}
Then the exact block energies satisfy
\[
\frac{d}{dt}\sum_j \alpha_j E_j(t)\ \le\ -\,2\,(1-\theta)\,\Gamma_{\min}\,\sum_j \alpha_j E_j(t),\qquad
\theta:=\sup_{t\in I}\overline{\Theta}(t)<1,
\]
for admissible positive weights $\alpha$ (equivalently the weighted norm $\|\cdot\|_{w,\op}$). Hence $\sum_j \alpha_j E_j$ decays exponentially on $I$.
\end{theorem}

\begin{proof}
(H1) places the PDE state in $B(t)$; the flux map is Lipschitz on bounded sets. Directed rounding in \eqref{eq:op-bound} encloses the true $\|\mathcal F(t)\|_{w,\op}$, so $\Theta(t)\le\overline{\Theta}(t)$. Apply the analytic small-gain lemma (Section~\ref{sec:analytic-core}) with $\theta=\sup\overline{\Theta}(t)$.
\end{proof}

\subsection{Archival and reproducibility}\label{subsec:cav-archival}
Artifacts for third-party re-verification:
\begin{itemize}
\item \textbf{Setup.} Grid $N$, viscosity $\nu$, $\Delta t$, domain, dealiasing mask, block partition $\{S_j\}$.
\item \textbf{Determinism.} RNG seeds, SHA-256 of masks/partitions and code snapshot, compiler flags.
\item \textbf{Certificates.} Per-step Krawczyk radii, FFT budgets ($K_{\mathrm{fft}}$, $\gamma_c(N)$), interval bounds for $\|\,|W^{1/2}\mathcal F W^{-1/2}|\,\|_{1,\infty}$, resulting $\overline{\Theta}(t)$, and weights $W$.
\item \textbf{Data format.} Enclosures as rational endpoints or fixed-precision decimals with declared endianness \cite{Tucker2011}.
\end{itemize}

\medskip
\noindent\textbf{Remark.} The computer-assisted route is conditional: it certifies that the small-gain hypothesis holds for the \emph{exact} PDE on the validated window, thereby closing the bootstrap in Sections~\ref{sec:analytic-core}–\ref{sec:bootstrap}.

\section{Extensions and stability}\label{sec:extensions}

\subsection{External forcing and domains with boundaries}\label{subsec:forcing-boundary}

\paragraph{Forcing classes.}
Let $f\in L^2_{\mathrm{loc}}([0,\infty);H^{-1})$ (or $f\in L^1_{\mathrm{loc}}([0,\infty);L^2)$ with $\int_0^\infty\!\|f(t)\|_{H^{-1}}^2\,dt<\infty$). The global energy balance is
\[
\frac{d}{dt}\tfrac12\|u\|_{L^2}^2+\nu\|\nabla u\|_{L^2}^2=\langle f,u\rangle.
\]
On blocks $S_j$,
\[
\frac{dE_j}{dt}
\;\le\; -\,2\,\Gamma_j E_j\;+\;\sum_i \mathcal F_{ji}\;+\;\langle f,u\rangle_{S_j},
\qquad \Gamma_j:=\nu\min_{k\in S_j}|k|^2,
\]
where the factor $2$ comes from $E_j=\tfrac12\sum_{k\in S_j}|\widehat u(k)|^2$. With weights $\alpha_j>0$,
\[
\frac{d}{dt}\mathcal E_\alpha
\;\le\; -\,2(1-\theta)\Gamma_{\min}\,\mathcal E_\alpha
\;+\;\sum_j \alpha_j\langle f,u\rangle_{S_j},
\qquad \mathcal E_\alpha:=\sum_j \alpha_j E_j,
\]
whenever $\|\mathcal F(t)\|_{w,\op}\le \theta\,\Gamma_{\min}$ with $0<\theta<1$.
By Cauchy–Schwarz and Young, for any $\varepsilon\in\big(0,2(1-\theta)\Gamma_{\min}\big)$,
\[
\sum_j \alpha_j\langle f,u\rangle_{S_j}
\;\le\; \frac{C_\alpha}{\varepsilon}\,\|f\|_{H^{-1}}^2 \;+\; \varepsilon\,\sum_j \alpha_j E_j,
\]
for some $C_\alpha$ depending on $\alpha$ and the block geometry. Hence
\[
\frac{d}{dt}\mathcal E_\alpha \;\le\; -\big(2(1-\theta)\Gamma_{\min}-\varepsilon\big)\,\mathcal E_\alpha
\;+\;\frac{C_\alpha}{\varepsilon}\,\|f\|_{H^{-1}}^2,
\]
which yields an absorbing bound and, for stationary/tempered forcing, uniform-in-time control (in 2D: compact global attractors) \cite{Temam2001,ConstantinFoias1988}.

\paragraph{Boundaries.}
On a bounded $C^{1,1}$ domain $\Omega\subset\R^3$ with no-slip $u|_{\partial\Omega}=0$, replace Fourier modes by the Stokes eigenbasis:
\[
-\Delta_S\phi_k=\lambda_k\phi_k\quad\text{in }L^2_\sigma(\Omega),\qquad 0<\lambda_1\le\lambda_2\le\cdots.
\]
Define spectral blocks $S_j=\{k:2^{2j}\le \lambda_k<2^{2(j+1)}\}$ and set $\Gamma_j=\nu\min_{k\in S_j}\lambda_k$, $\Gamma_{\min}\ge\nu\lambda_1(\Omega)$. The Stokes semigroup is analytic on $L^p_\sigma(\Omega)$ for $1<p<\infty$ \cite{Sohr2001,Temam2001}. The trilinear form satisfies $b(u,v,v)=0$ for $u,v\in H^1_0$ with $\nabla\!\cdot u=0$, so the rotational/Leray form remains $L^2$–skew and $\mathcal F_{ij}=-\mathcal F_{ji}$ persists. Spectral Galerkin de-aliasing is handled by exact projection/bandwidth separation (no FFT $2/3$ rule). The certificate reads $\|\mathcal F\|_{w,\op}\le \theta\,\Gamma_{\min}$ with rates depending on $\lambda_1(\Omega)$.

\subsection{Endpoint/near-critical spaces and rough data}\label{subsec:endpoint-rough}
Small-data global well-posedness holds in the critical classes $L^3(\R^3)$, $\dot H^{1/2}$, $\mathrm{BMO}^{-1}$, and $\dot B^{-1+3/p}_{p,\infty}$ \cite{Kato1984,Cannone1995,KochTataru2001,Lemarie2002}. Choosing
\[
\alpha_j\asymp 2^{2\sigma j}\quad(\sigma>1/2)\qquad\Longrightarrow\qquad
\mathcal E_\alpha(t)\asymp \|u(t)\|_{\dot H^\sigma}^2=:X_\sigma(t),
\]
the small-gain inequality yields $X_\sigma(t)\le X_\sigma(0)e^{-2(1-\theta)\Gamma_{\min}t}$ and, by $\dot H^\sigma\hookrightarrow L^q$ with $q=\frac{6}{3-2\sigma}>3$, a Serrin bound $u\in L^p_tL^q_x$ with $2/p+3/q=1$. This gives continuation and uniqueness in the strong class.
For large rough data, decompose $u_0=u_{0,\mathrm{lo}}+u_{0,\mathrm{hi}}$ with $u_{0,\mathrm{hi}}$ small in a critical norm; propagate $u_{\mathrm{hi}}$ by small-data theory and control $u_{\mathrm{lo}}$ on a short interval via small-gain, then iterate. Stability under perturbations follows from the semigroup Lipschitz dependence in the controlled norms and the antisymmetry of $\mathcal F$.

\subsection{2D vorticity check and Euler limits $\nu\to0$}\label{subsec:2D-euler}
In 2D,
\[
\partial_t\omega+u\cdot\nabla\omega=\nu\Delta\omega+g,\qquad
\|\omega(t)\|_{L^\infty}\le \|\omega_0\|_{L^\infty}+\int_0^t\|g(s)\|_{L^\infty}\,ds,
\]
and $\tfrac12\frac{d}{dt}\|\omega\|_{L^2}^2+\nu\|\nabla\omega\|_{L^2}^2=\langle g,\omega\rangle$ \cite{Ladyzhenskaya1969,Temam2001,MajdaBertozzi2002}. With $\Gamma_{\min}=\nu k_{\min}^2$ on $\T^2$, the certificate gives exponential decay of $\dot H^\sigma(\T^2)$, consistent with classical theory. For Euler on $\T^2$ with $\omega_0\in L^\infty$, Yudovich ensures global well-posedness \cite{Yudovich1963}. As $\nu\to0$, on periodic domains and for smooth data, $u^\nu\to u^{\mathrm{Euler}}$ on any Euler existence interval \cite{MajdaBertozzi2002}; with boundaries, Kato’s criterion characterizes convergence via vanishing viscous dissipation in an $O(\nu)$ layer \cite{KatoZero1984}. The small-gain framework transfers to the Stokes spectrum in 2D domains and matches these limits.

\section{Discussion of unavoidable gaps and how they are sealed}\label{sec:gaps}

\subsection{Classical inputs and exact points of invocation}\label{subsec:classical-inputs}
We use the following standard results at the indicated steps:
\begin{itemize}
  \item \emph{Leray–Hopf existence and energy inequality} on $\T^3$ or $\R^3$ for $u_0\in L^2_\sigma$, $f\in L^2_{\mathrm{loc}}(\R_+;H^{-1})$: construction and a priori bounds in \S\ref{sec:leray-hopf}. \cite{Leray1934,Hopf1951,Temam2001,Sohr2001}
  \item \emph{Local energy inequality / suitable weak solutions} for the partial-regularity scaffold in \S\ref{sec:ckn}. \cite{CKN1982,Lin1998}
  \item \emph{Local strong/mild well-posedness in critical spaces} (Fujita–Kato; Koch–Tataru) and \emph{continuation criteria} (Serrin; Beale–Kato–Majda), invoked in \S\ref{sec:local} and \S\ref{sec:bootstrap}. \cite{FujitaKato1964,Kato1984,KochTataru2001,Serrin1962,BealeKatoMajda1984,EscauriazaSereginSverak2003}
  \item \emph{Pressure estimates via Riesz transforms} (Calderón–Zygmund), used in \S\ref{sec:apriori}. \cite{Stein1993,Temam2001}
  \item \emph{2D sanity checks} (Yudovich) and classical decay, used in \S\ref{subsec:2D-euler}. \cite{Yudovich1963,MajdaBertozzi2002}
\end{itemize}

\subsection{New constants and reproducible derivations}\label{subsec:new-constants}
All constants from the block-energy method and the validation layer are explicit and auditable:
\begin{enumerate}
  \item \textbf{Block dissipation floor.} For $S_j$,
  \[
  \frac{dE_j}{dt}\Big|_{\mathrm{lin}}
  =-\nu\!\sum_{k\in S_j}\!|k|^2|\widehat u(k)|^2
  \le -2\,\Gamma_j E_j,\qquad
  \Gamma_j:=\nu\min_{k\in S_j}|k|^2,
  \]
  since $E_j=\tfrac12\sum_{k\in S_j}|\widehat u(k)|^2$ (\S\ref{sec:analytic-core}).
  \item \textbf{Flux bound and small-gain constant.} With $\mathcal F^\top=-\mathcal F$ and weights $\alpha_j\simeq w_j>0$,
  \[
  \sum_{i,j}\alpha_j\mathcal F_{ji}
  =\tfrac12\sum_{i,j}(\alpha_j-\alpha_i)\mathcal F_{ji}
  \ \ \Rightarrow\ \
  \Big|\sum_{i,j}\alpha_j\mathcal F_{ji}\Big|
  \le \|W^{1/2}\mathcal F W^{-1/2}\|_2 \sum_j \alpha_j E_j,
  \]
  with $W=\mathrm{diag}(w_j)$. We certify $\|W^{1/2}\mathcal F W^{-1/2}\|_2$ by
  \[
  \|A\|_2\ \le\ \sqrt{\ \|\,|A|\,\|_1\ \|\,|A|\,\|_\infty\ },
  \]
  applied to $A=W^{1/2}\mathcal F W^{-1/2}$ (\S\ref{sec:cav}).
  \item \textbf{Rotational-kernel magnitude.} For triads $k+\ell+m=0$ in rotational form,
  \[
  \big|\overline{\widehat u(k)}\cdot\big[(\ell\times \widehat u(\ell))\times \widehat u(m)\big]\big|
  \ \le\ (|k|+|\ell|)\,|\widehat u(k)|\,|\widehat u(\ell)|\,|\widehat u(m)|,
  \]
  yielding an assembly of $|\mathcal F|$ from modewise magnitudes (no cancellation).
  \item \textbf{Semigroup/interpolation constants.} Heat/Stokes smoothing, Sobolev, and Gagliardo–Nirenberg/Ladyzhenskaya constants are taken in standard forms on $\T^3$ (or via the Stokes spectrum on bounded domains). \cite{Temam2001,Sohr2001}
  \item \textbf{Validated numerics constants.} (i) FFT roundoff via Higham’s $\gamma_c(N)=\dfrac{c\log_2N\,\varepsilon}{1-c\log_2N\,\varepsilon}$; (ii) Krawczyk one-step enclosure with SSP–RK3 local defect and semigroup-based Lipschitz radii; (iii) interval evaluation of $\|\,|A|\,\|_{1,\infty}$. \cite{Higham2002,Neumaier1990,Rump1999,Tucker2011}
\end{enumerate}
Each constant is emitted by scripts (block dissipation, flux assembly, norm ledgers, and validation logs) with fixed seeds and code hashes (\S\ref{sec:cav}).

\subsection{Limitations and conditional hypotheses}\label{subsec:limitations}
\begin{itemize}
  \item \textbf{Certificate hypothesis.} The key bound $\|\mathcal F\|_{w,\op}\le \theta\,\Gamma_{\min}$ with $0<\theta<1$ is either proven analytically for all $t$ (unconditional conclusions) or verified on finite windows (\S\ref{sec:cav}); in the latter case results are conditional on covering time by validated windows.
  \item \textbf{Discretization dependence (validated path).} The certified route assumes a structure-preserving scheme (Leray projection $+$ rotational form $+$ de-aliasing). Other discretizations may not yield equally tight bounds on $\mathcal F$.
  \item \textbf{Endpoint/rough-data regimes.} The bootstrap closes once a Serrin-type bound is secured (\S\ref{sec:bootstrap}). Endpoints and large critical data rely on \cite{KochTataru2001,EscauriazaSereginSverak2003} and remain delicate for uniform certification.
  \item \textbf{Boundaries and $\nu\to0$.} With walls, dissipation floors involve the Stokes spectrum and boundary layers; Kato’s criterion governs the Euler limit \cite{KatoZero1984}. Our framework adapts spectrally but sharp zero-viscosity limits near walls lie beyond scope.
\end{itemize}

\medskip
\noindent\textbf{Reproducibility.} All nonstandard constants/certificates are produced by scripts emitting machine-independent manifests: RNG seeds, masks/partitions, code/version hashes, and interval/norm ledgers (\S\ref{sec:cav}). This seals computational gaps to a standard of rigorous a~posteriori verification while keeping the analytic chain explicit.

\appendix
\section{Preliminaries and technical lemmas}\label{sec:appendix}

\subsection*{A. Calderón–Zygmund, Riesz boundedness, and pressure formulae}\label{app:CZ}
\begin{lemma}[Riesz transform bounds]\label{lem:Riesz}
For $1<p<\infty$ the Riesz transforms $R_j=\partial_j(-\Delta)^{-1/2}$ are bounded on $L^p(\R^3)$ and on $L^p(\T^3)$ (with zero mean). In particular,
\[
\|R_j f\|_{L^p}\le C_p\|f\|_{L^p},\qquad
\|\nabla(-\Delta)^{-1} g\|_{W^{1,p}}\le C_p\|g\|_{L^p}.
\]
\end{lemma}
\begin{proof}[Sketch]
Classical Calderón–Zygmund theory; periodic case via Fourier multipliers. \cite[Ch.~IV]{Stein1993}.
\end{proof}

\begin{lemma}[Pressure representation]\label{lem:pressure}
Let $(u,p)$ solve NSE with $\nabla\!\cdot u=0$ on $\Omega\in\{\R^3,\T^3\}$. Then
\[
-\Delta p=\partial_i\partial_j(u_i u_j)-\nabla\!\cdot f,\qquad
p=\sum_{i,j=1}^3 R_iR_j(u_i u_j)-(-\Delta)^{-1}\nabla\!\cdot f,
\]
and for $1<p<\infty$,
\[
\|\nabla p\|_{L^p}\le C_p\|u\otimes u\|_{L^p}+C_p\|f\|_{L^p}.
\]
\end{lemma}
\begin{proof}[Sketch]
Take divergence, solve Poisson equation, apply Lemma~\ref{lem:Riesz}. \cite[§I.3, §III.3]{Temam2001}, \cite[Thm.~2.3.1]{Sohr2001}.
\end{proof}

\subsection*{B. Semigroup bounds and Duhamel estimates}\label{app:SG}
\begin{lemma}[Heat/Stokes semigroup]\label{lem:heat}
For $t>0$, $1\le p\le q\le\infty$, $m\in\N$,
\[
\|e^{\nu t\Delta} f\|_{L^q}\le C\,(\nu t)^{-\frac{3}{2}\left(\frac1p-\frac1q\right)}\|f\|_{L^p},\qquad
\|\nabla^m e^{\nu t\Delta} f\|_{L^p}\le C\,(\nu t)^{-m/2}\|f\|_{L^p},
\]
and the analogous bounds hold for the Stokes semigroup on $L^p_\sigma$. 
\end{lemma}
\begin{proof}[Sketch]
Heat kernel estimates; periodic case by Fourier series; Stokes via analyticity. \cite[Ch.~3]{Temam2001}, \cite[Thm.~2.6]{Sohr2001}.
\end{proof}

\begin{lemma}[Mild formulation and bilinear estimate]\label{lem:duhamel}
For divergence-free $u$,
\[
u(t)=e^{\nu t\Delta}u_0-\int_0^t e^{\nu(t-s)\Delta}\mathbb P\nabla\!\cdot(u\otimes u)(s)\,ds+\int_0^t e^{\nu(t-s)\Delta}\mathbb P f(s)\,ds.
\]
If $X$ is a critical space (e.g.\ $L^3$, $\dot H^{1/2}$, $BMO^{-1}$), then
\[
\Big\|\int_0^t e^{\nu(t-s)\Delta}\mathbb P\nabla\!\cdot(u\otimes v)(s)\,ds\Big\|_{X}
\le C\,\|u\|_{X}\|v\|_{X}.
\]
\end{lemma}
\begin{proof}[Sketch]
Semigroup mapping properties and product estimates defining $X$. \cite{FujitaKato1964,Kato1984,KochTataru2001}.
\end{proof}

\subsection*{C. $\varepsilon$–regularity statement}\label{app:eps}
\begin{lemma}[CKN $\varepsilon$–regularity]\label{lem:ckn}
There exist universal $\varepsilon_0,\beta>0$ such that if $(u,p)$ is suitable in $Q_1=B_1\times(-1,0)$ with
\[
\iint_{Q_1}\big(|u|^3+|p|^{3/2}\big)\,dx\,dt\le \varepsilon_0,
\]
then $u$ is Hölder continuous in $Q_{1/2}$ and $\|u\|_{C^\beta(Q_{1/2})}\le C$.
\end{lemma}
\begin{proof}[Sketch]
CKN local energy method (and Lin’s simplification). \cite{CKN1982,Lin1998}.
\end{proof}

\subsection*{D. Discrete skew-symmetry and Leray commutation}\label{app:discrete}
\begin{lemma}[Skew-symmetry of the rotational form]\label{lem:skew}
Implement the nonlinearity as $\widehat{\omega\times u}$ with $\omega=\nabla\times u$, apply $\mathbb P$ modewise, and use $2/3$ de-aliasing. For inviscid, unforced dynamics,
\[
\sum_k \widehat u(k)\cdot \overline{\widehat{(\omega\times u)}(k)}\in i\R,
\]
so the discrete kinetic energy is preserved up to roundoff at each SSP–RK3 stage.
\end{lemma}
\begin{proof}[Sketch]
Triadic antisymmetry of the rotational kernel and exact triad closure under $2/3$-rule. \cite[§3.6]{CanutoQuarteroniHussainiZang2006}, \cite{PattersonOrszag1971,Morinishi1998}.
\end{proof}

\begin{lemma}[Commutation with spectral truncation]\label{lem:commute}
If $K_N$ is symmetric and $\Pi_{K_N}$ the Fourier projector, then $\mathbb P\Pi_{K_N}=\Pi_{K_N}\mathbb P$.
\end{lemma}
\begin{proof}[Sketch]
$\mathbb P$ is a diagonal Fourier multiplier $(I-\frac{k\otimes k}{|k|^2})$ preserving $K_N$. \cite[§I.1]{Temam2001}.
\end{proof}

\subsection*{E. Interval arithmetic and FFT-error lemmas}\label{app:interval}
\begin{lemma}[Operator-norm enclosure]\label{lem:opnorm}
For any real matrix $A$,
\[
\|A\|_2\le \sqrt{\|\,|A|\,\|_1\,\|\,|A|\,\|_\infty}.
\]
If $|A|\le B$ entrywise, then $\|A\|_2\le \sqrt{\|B\|_1\|B\|_\infty}$.
\end{lemma}
\begin{proof}
\cite[Thm.~5.6.18]{HornJohnsonMatrix}.
\end{proof}

\begin{lemma}[FFT roundoff accumulation]\label{lem:fft}
Let $\widehat x$ be the exact FFT of $x\in\C^{N\times N}$ and $\widetilde{\widehat x}$ the computed one in IEEE double. Then
\[
\|\widehat x-\widetilde{\widehat x}\|_{\ell^\infty}\le \gamma_c(N)\,\|x\|_{\ell^1},\qquad
\gamma_c(N)=\frac{10\,(\log_2 N)\,u}{1-10\,(\log_2 N)\,u},
\]
with unit roundoff $u=2^{-53}$. For $K$ transforms in a step, the cumulative bound is $K\,\gamma_c(N)\,\|x\|_{\ell^1}$; using $\|x\|_{\ell^1}\le \sqrt{N^2}\|x\|_{\ell^2}$ yields a computable enclosure.
\end{lemma}
\begin{proof}[Sketch]
Higham’s rounding-error model for FFTs with a conservative operation count. \cite[Ch.~22]{Higham2002}.
\end{proof}

\subsection*{F. Krawczyk operator and local truncation constants}\label{app:krawczyk}
\begin{lemma}[Krawczyk inclusion for one time step]\label{lem:krawczyk}
Let $F(u)=u-\Phi_{\Delta t}(u)$ for the projected SSP–RK3 map $\Phi_{\Delta t}$. Choose $M\approx F'(u_c)^{-1}$ diagonal in Fourier (linear part) and a ball $B$ around a predictor $u_c$. If
\[
q:=\|M\|\,\Delta t\,L(u_c)<1,\qquad
r:= C_4\,\Delta t^4\,L(u_c)^4\,\|u_c\|_{L^2}\quad (C_4=1/24),
\]
with $L(u_c)\le \nu k_{\max}^2+2\|\nabla u_c\|_{L^\infty}$, then the exact solution at $t+\Delta t$ lies in $B\big(u_{\mathrm{num}},\, r/(1-q)\big)$ about the numerical value $u_{\mathrm{num}}$.
\end{lemma}
\begin{proof}[Sketch]
Standard Krawczyk/inclusion argument with SSP–RK3 local defect and bilinear Lipschitz control. \cite{Neumaier1990,Rump1999,Tucker2011,ButcherBook}.
\end{proof}

\medskip
\noindent\textbf{Cross-references.}
Lemmas~\ref{lem:Riesz}–\ref{lem:pressure} feed \S\ref{sec:apriori};
Lemmas~\ref{lem:heat}–\ref{lem:duhamel} underlie \S\ref{sec:local};
Lemma~\ref{lem:ckn} underpins \S\ref{sec:ckn};
Lemmas~\ref{lem:skew}–\ref{lem:commute} justify \S\ref{sec:cav};
Lemmas~\ref{lem:opnorm}–\ref{lem:fft} and \ref{lem:krawczyk} establish the numerical layer in \S\ref{sec:cav}.

% Bibliography keys used above:
% Stein1993; Temam2001; Sohr2001; FujitaKato1964; Kato1984; KochTataru2001; CKN1982; Lin1998;
% CanutoQuarteroniHussainiZang2006; PattersonOrszag1971; Morinishi1998; Higham2002;
% HornJohnsonMatrix; Neumaier1990; Rump1999; Tucker2011; ButcherBook.



\section*{Acknowledgments}
% optional

% --- NYJM bibliography model: inline thebibliography with MR/Zbl/DOI/arXiv ---
% Example entries; replace with actual items and add \mrev{} and \zbl{} macros if you use them.


\begin{thebibliography}{99}


\begin{thebibliography}{99}

\bibitem{BahouriCheminDanchin2011}
H.~Bahouri, J.-Y.~Chemin, R.~Danchin,
\emph{Fourier Analysis and Nonlinear Partial Differential Equations},
Grundlehren der Math. Wissenschaften \textbf{343}, Springer, 2011.
% \mrev{} \zbl{} \doi{}

\bibitem{BahouriGerard1999}
H.~Bahouri, P.~Gérard,
High frequency approximation of solutions to critical nonlinear wave equations,
\emph{Amer. J. Math.} \textbf{121} (1999), 131--175.
% \mrev{} \zbl{} \doi{}

\bibitem{BealeKatoMajda1984}
J.~T.~Beale, T.~Kato, A.~Majda,
Remarks on the breakdown of smooth solutions for the 3-D Euler equations,
\emph{Comm. Math. Phys.} \textbf{94} (1984), 61--66.
% \mrev{} \zbl{} \doi{}

\bibitem{Boyd2001}
J.~P.~Boyd,
\emph{Chebyshev and Fourier Spectral Methods}, 2nd ed.,
Dover, 2001.
% \mrev{} \zbl{}

\bibitem{ButcherBook}
J.~C.~Butcher,
\emph{Numerical Methods for Ordinary Differential Equations}, 2nd ed.,
Wiley, 2008.
% \mrev{} \zbl{}

\bibitem{Cannone1995}
M.~Cannone,
\emph{Ondelettes, paraproduits et Navier--Stokes},
Diderot, Paris, 1995.
% \mrev{} \zbl{}

\bibitem{CannoneBook}
M.~Cannone,
\emph{Harmonic Analysis Tools for the Navier--Stokes Equations},
AMS, 2004.
% \mrev{} \zbl{}

\bibitem{CKN1982}
L.~Caffarelli, R.~Kohn, L.~Nirenberg,
Partial regularity of suitable weak solutions of the Navier--Stokes equations,
\emph{Comm. Pure Appl. Math.} \textbf{35} (1982), 771--831.
% \mrev{} \zbl{} \doi{}

\bibitem{CanutoQuarteroniHussainiZang2006}
C.~Canuto, M.~Y.~Hussaini, A.~Quarteroni, T.~A.~Zang,
\emph{Spectral Methods: Fundamentals in Single Domains},
Springer, 2006.
% \mrev{} \zbl{}

\bibitem{ConstantinFoias1988}
P.~Constantin, C.~Foias,
\emph{Navier--Stokes Equations},
Univ. of Chicago Press, 1988.
% \mrev{} \zbl{}

\bibitem{EscauriazaSereginSverak2003}
L.~Escauriaza, G.~Seregin, V.~\v{S}ver\'ak,
$L^{3,\infty}$-solutions of Navier--Stokes and backward uniqueness,
\emph{Uspekhi Mat. Nauk} \textbf{58} (2003), 3--44.
% \mrev{} \zbl{} \doi{}

\bibitem{FujitaKato1964}
H.~Fujita, T.~Kato,
On the Navier--Stokes initial value problem I,
\emph{Arch. Ration. Mech. Anal.} \textbf{16} (1964), 269--315.
% \mrev{} \zbl{}

\bibitem{Frisch1995}
U.~Frisch,
\emph{Turbulence: The Legacy of A. N. Kolmogorov},
Cambridge Univ. Press, 1995.
% \mrev{} \zbl{}

\bibitem{GallagherKochPlanchon2016}
I.~Gallagher, H.~Koch, F.~Planchon,
A profile decomposition approach to the $L^3$ Navier--Stokes regularity problem,
\emph{ESAIM Proc. Surveys} \textbf{48} (2015), 1--20.
% \mrev{} \zbl{}

\bibitem{Gagliardo1959}
E.~Gagliardo,
Ulteriori propriet\`a di alcune classi di funzioni in pi\`u variabili,
\emph{Ricerche Mat.} \textbf{8} (1959), 24--51.
% \mrev{} \zbl{}

\bibitem{Gerard1998}
P.~G\'erard,
Description du d\'efaut de compacit\'e de l’injection de Sobolev,
\emph{ESAIM: COCV} \textbf{3} (1998), 213--233.
% \mrev{} \zbl{}

\bibitem{GigaMild1986}
Y.~Giga,
Solutions for semilinear parabolic equations in $L^p$ and regularity of weak solutions of the Navier--Stokes system,
\emph{J. Differential Equations} \textbf{62} (1986), 186--212.
% \mrev{} \zbl{}

\bibitem{Giga1981}
Y.~Giga,
Analyticity of the semigroup generated by the Stokes operator in $L^r$ spaces,
\emph{Math. Z.} \textbf{178} (1981), 297--329.
% \mrev{} \zbl{}

\bibitem{Higham2002}
N.~J.~Higham,
\emph{Accuracy and Stability of Numerical Algorithms}, 2nd ed.,
SIAM, 2002.
% \mrev{} \zbl{}

\bibitem{HornJohnsonMatrix}
R.~A.~Horn, C.~R.~Johnson,
\emph{Matrix Analysis}, 2nd ed.,
Cambridge Univ. Press, 2013.
% \mrev{} \zbl{}

\bibitem{Hopf1951}
E.~Hopf,
\"Uber die Anfangswertaufgabe f\"ur die hydrodynamischen Grundgleichungen,
\emph{Math. Nachr.} \textbf{4} (1951), 213--231.
% \mrev{} \zbl{}

\bibitem{JiangTeelPraly1994}
Z.-P.~Jiang, A.~R.~Teel, L.~Praly,
Small-gain theorem for ISS systems,
\emph{Math. Control Signals Systems} \textbf{7} (1994), 95--120.
% \mrev{} \zbl{}

\bibitem{Kato1984}
T.~Kato,
Strong $L^p$-solutions of the Navier--Stokes equation in $\R^m$,
\emph{Math. Z.} \textbf{187} (1984), 471--480.
% \mrev{} \zbl{}

\bibitem{KatoZero1984}
T.~Kato,
Remarks on zero viscosity limit for incompressible flow with boundary,
in \emph{Seminar on Nonlinear PDEs}, MSRI Publ. \textbf{2}, 1984, 85--98.
% \mrev{} \zbl{}

\bibitem{KenigMerle2006}
C.~E.~Kenig, F.~Merle,
Global well-posedness, scattering and blow-up for the energy-critical focusing NLS,
\emph{Acta Math.} \textbf{201} (2008), 147--212.
% \mrev{} \zbl{}

\bibitem{Khalil2002}
H.~K.~Khalil,
\emph{Nonlinear Systems}, 3rd ed.,
Prentice Hall, 2002.
% \mrev{} \zbl{}

\bibitem{KochTataru2001}
H.~Koch, D.~Tataru,
Well-posedness for the Navier--Stokes equations,
\emph{Adv. Math.} \textbf{157} (2001), 22--35.
% \mrev{} \zbl{} \doi{}

\bibitem{KozonoTaniuchi2000}
H.~Kozono, Y.~Taniuchi,
Bilinear estimates in BMO and the Navier--Stokes equations,
\emph{Math. Z.} \textbf{235} (2000), 173--194.
% \mrev{} \zbl{}

\bibitem{Ladyzhenskaya1969}
O.~A.~Ladyzhenskaya,
\emph{The Mathematical Theory of Viscous Incompressible Flow}, 2nd ed.,
Gordon \& Breach, 1969.
% \mrev{} \zbl{}

\bibitem{LadyzhenskayaSereginSverak1999}
O.~A.~Ladyzhenskaya, G.~Seregin, V.~\v{S}ver\'ak,
Weak solutions with initial data in $L^p$,
\emph{J. Math. Fluid Mech.} \textbf{1} (1999), 356--364.
% \mrev{} \zbl{}

\bibitem{Lemarie2002}
P.-G.~Lemari\'e--Rieusset,
\emph{Recent Developments in the Navier--Stokes Problem},
CRC Press, 2002.
% \mrev{} \zbl{}

\bibitem{Leray1934}
J.~Leray,
Sur le mouvement d’un liquide visqueux emplissant l’espace,
\emph{Acta Math.} \textbf{63} (1934), 193--248.
% \mrev{} \zbl{}

\bibitem{Lin1998}
F.-H.~Lin,
A new proof of the Caffarelli--Kohn--Nirenberg theorem,
\emph{Comm. Pure Appl. Math.} \textbf{51} (1998), 241--257.
% \mrev{} \zbl{}

\bibitem{Lions1969}
J.-L.~Lions,
\emph{Quelques m\'ethodes de r\'esolution des probl\`emes aux limites non lin\'eaires},
Dunod; Gauthier-Villars, Paris, 1969.
% \mrev{} \zbl{}

\bibitem{MajdaBertozzi2002}
A.~J.~Majda, A.~L.~Bertozzi,
\emph{Vorticity and Incompressible Flow},
Cambridge Univ. Press, 2002.
% \mrev{} \zbl{}

\bibitem{Morinishi1998}
Y.~Morinishi, T.~Lund, O.~Vasilyev, P.~Moin,
Fully conservative higher order finite difference schemes for incompressible flow,
\emph{J. Comput. Phys.} \textbf{143} (1998), 90--124.
% \mrev{} \zbl{}

\bibitem{Moore1966}
R.~E.~Moore,
\emph{Interval Analysis},
Prentice Hall, 1966.
% \mrev{} \zbl{}

\bibitem{Neumaier1990}
A.~Neumaier,
\emph{Interval Methods for Systems of Equations},
Cambridge Univ. Press, 1990.
% \mrev{} \zbl{}

\bibitem{NecasRuzickaSverak1996}
J.~Ne\v{c}as, M.~R{\r u}\v{z}i\v{c}ka, V.~\v{S}ver\'ak,
On Leray’s self-similar solutions,
\emph{Acta Math.} \textbf{176} (1996), 283--294.
% \mrev{} \zbl{}

\bibitem{Nirenberg1959}
L.~Nirenberg,
On elliptic partial differential equations,
\emph{Ann. Sc. Norm. Sup. Pisa} \textbf{13} (1959), 115--162.
% \mrev{} \zbl{}

\bibitem{Orszag1971}
S.~A.~Orszag,
On the elimination of aliasing in finite-difference schemes by filtering high-wavenumber components,
\emph{J. Atmos. Sci.} \textbf{28} (1971), 1074.
% \mrev{} \zbl{}

\bibitem{PattersonOrszag1971}
G.~S.~Patterson, S.~A.~Orszag,
Spectral calculations of isotropic turbulence: Efficient removal of aliasing interactions,
\emph{Phys. Fluids} \textbf{14} (1971), 2538--2541.
% \mrev{} \zbl{}

\bibitem{Prodi1959}
G.~Prodi,
Un teorema di unicità per le equazioni di Navier--Stokes,
\emph{Ann. Mat. Pura Appl.} \textbf{48} (1959), 173--182.
% \mrev{} \zbl{}

\bibitem{Rump1999}
S.~M.~Rump,
INTLAB---Interval laboratory,
in \emph{Developments in Reliable Computing}, Kluwer, 1999, 77--104.
% \mrev{} \zbl{}

\bibitem{Serrin1962}
J.~Serrin,
On the interior regularity of weak solutions of the Navier--Stokes equations,
\emph{Arch. Ration. Mech. Anal.} \textbf{9} (1962), 187--195.
% \mrev{} \zbl{}

\bibitem{Serrin1963}
J.~Serrin,
The initial value problem for the Navier--Stokes equations,
in \emph{Nonlinear Problems} (Proc. Sympos., Madison, 1962),
Univ. of Wisconsin Press, 1963, 69--98.
% \mrev{} \zbl{}

\bibitem{Sohr2001}
H.~Sohr,
\emph{The Navier--Stokes Equations: An Elementary Functional Analytic Approach},
Birkh\"auser, 2001.
% \mrev{} \zbl{}

\bibitem{Stein1993}
E.~M.~Stein,
\emph{Harmonic Analysis: Real-Variable Methods, Orthogonality, and Oscillatory Integrals},
Princeton Univ. Press, 1993.
% \mrev{} \zbl{}

\bibitem{Temam2001}
R.~Temam,
\emph{Navier--Stokes Equations: Theory and Numerical Analysis},
AMS Chelsea, 2001.
% \mrev{} \zbl{}

\bibitem{Tsai1998}
T.-P.~Tsai,
On Leray’s self-similar solutions of the Navier--Stokes equations,
\emph{Comm. Partial Differential Equations} \textbf{27} (2002), 2363--2379.
% \mrev{} \zbl{}

\bibitem{Tucker2011}
W.~Tucker,
\emph{Validated Numerics},
Princeton Univ. Press, 2011.
% \mrev{} \zbl{}

\bibitem{Vasseur2007}
A.~Vasseur,
A new proof of partial regularity of solutions to Navier--Stokes equations,
\emph{NoDEA} \textbf{14} (2007), 753--785.
% \mrev{} \zbl{}

\bibitem{Yudovich1963}
V.~I.~Yudovich,
Non-stationary flows of an ideal incompressible fluid,
\emph{Z. Vych. Mat.} \textbf{3} (1963), 1032--1066.
% \mrev{} \zbl{}

\end{thebibliography}


\end{document}
